\documentclass[12pt, titlepage]{article}

\usepackage{amsmath, mathtools}

\usepackage[round]{natbib}
\usepackage{amsfonts}
\usepackage{amssymb}
\usepackage{graphicx}
\usepackage{colortbl}
\usepackage{xr}
\usepackage{hyperref}
\usepackage{longtable}
\usepackage{xfrac}
\usepackage{tabularx}
\usepackage{float}
\usepackage{enumitem}

\usepackage{booktabs}
\usepackage{multirow}
\usepackage[section]{placeins}
\usepackage{caption}
\usepackage{fullpage}
\usepackage[table]{xcolor}
\usepackage{array} 
\usepackage{colortbl}
\usepackage{makecell}
\usepackage{longtable}
\usepackage{centernot}
\usepackage{verbatim}
\usepackage[edges]{forest}

\hypersetup{
bookmarks=true,     % show bookmarks bar?
colorlinks=true,       % false: boxed links; true: colored links
linkcolor=red,          % color of internal links (change box color with linkbordercolor)
citecolor=blue,      % color of links to bibliography
filecolor=magenta,  % color of file links
urlcolor=cyan          % color of external links
}

\definecolor{tableblue}{RGB}{230, 240, 255}
\definecolor{headerblue}{RGB}{100, 150, 220}
\definecolor{accentmint}{RGB}{170,240,209} 
\definecolor{darkpurple}{RGB}{75,0,130}      % Royal purple
\definecolor{lavender}{RGB}{230,230,250}     % Lavender
\definecolor{lightlavender}{RGB}{245,245,255} % Light lavender
\definecolor{green}{RGB}{0,80,50}
\definecolor{lightgreen}{RGB}{0,200,120}

%new colours for final draft
\definecolor{headerblue}{RGB}{41, 128, 185}
\definecolor{lightblue}{RGB}{232, 244, 250}
\definecolor{softgray}{RGB}{245, 245, 245}
\definecolor{rowstripe}{RGB}{249, 249, 249}



  %\item Exercise data: this table consists of exercise instructions, maximum height (the highest point their leg can lift up to), minimum height (starting point) and user\_account\_p['Name'].
  %\item User activity: This will contain information about a specific patient's exercise activity. Data such as exercise duration, maximum height, number of repetitions, target area, the date the exercise was performed and its time, rest time, number of sets and analysis. Additionally, it will include a column for the patient's feedback on their exercise performance.


\newcolumntype{L}{>{\raggedright\arraybackslash}p{0.2\textwidth}}
\newcolumntype{M}{>{\raggedright\arraybackslash}p{0.15\textwidth}}
\newcolumntype{R}{>{\raggedright\arraybackslash}p{0.55\textwidth}}
\newcolumntype{S}{>{\raggedright\arraybackslash}p{0.26\textwidth}}


\externaldocument{../../SRS/SRS}

\input{../../Comments}
\input{../../Common}

\begin{document}

\title{Module Interface Specification for \progname{}}

\author{\authname}

\date{\today}

\maketitle

\pagenumbering{roman}

\section{Revision History}

\begin{tabularx}{\textwidth}{p{3cm}p{2cm}X}
\toprule {\bf Date} & {\bf Version} & {\bf Notes}\\
\midrule
Nov 6 - Nov 14 & 1.0 & Preliminary Draft\\
March 17 & 2.0 & Starting final draft\\
March 17 - March 19 & 2.0 & Started Account Creation module\\
March 30 - April 2 & 2.0 & Account Creation module work in progress\\
April 3 & 2.0 & Valid User module started\\
April 5 & 2.0 & User Account module complete\\
April 5 & 2.0 & Exercise Data module in progress\\
April 6 & 2.0 & Working on Analysis control flow module\\
April 7 & 2.0 & Final touches + database schema added\\
April 14 & 2.0 & User account, analyzer and feedback modules in progress\\




\bottomrule
\end{tabularx}
\\
\\

\section{Symbols, Abbreviations and Acronyms}

See SRS Documentation at \href{https://github.com/M9Huynh/technically-functional/blob/main/docs/SRS/SRS.pdf}{SRS}

\newpage

\tableofcontents

\newpage

\pagenumbering{arabic}

\section{Introduction}

%new
The following document details the Module Interface Specifications for the PhysioCompanion app. The application is 
designed to identify users performing prescribed exercises and provide feedback on their form. This tool uses computer
vision, specifically OpenCV and MediaPipe, to analyze user video recording and perform analysis. Please note that certain
modules exhibit functional overlap, as their respective operations perform equivalent decomposed tasks across different 
modules. \\

Complementary documents include the \href{https://github.com/M9Huynh/technically-functional/blob/main/docs/SRS/SRS.pdf}{System Requirement Specifications}
and \href{https://github.com/M9Huynh/technically-functional/blob/main/docs/Design/SoftArchitecture/MG.pdf}{Module Guide}.  The full documentation and implementation can be
found at \href{https://github.com/M9Huynh/technically-functional}{here}.  

\section{Notation}

The structure of the MIS for modules comes from \citet{HoffmanAndStrooper1995},
with the addition that template modules have been adapted from
\cite{GhezziEtAl2003}.  The mathematical notation comes from Chapter 3 of
\citet{HoffmanAndStrooper1995}.  For instance, the symbol := is used for a
multiple assignment statement and conditional rules follow the form $(c_1
\Rightarrow r_1 | c_2 \Rightarrow r_2 | ... | c_n \Rightarrow r_n )$.

The following table summarizes the primitive data types used by \progname. 

\begin{center}
\renewcommand{\arraystretch}{1.2}
\noindent 
\begin{tabular}{l l p{7.5cm}} 
\toprule 
\textbf{Data Type} & \textbf{Notation} & \textbf{Description}\\ 
\midrule
character & char & a single symbol or digit\\
integer & $\mathbb{Z}$ & a number without a fractional component in (-$\infty$, $\infty$) \\
natural number & $\mathbb{N}$ & a number without a fractional component in [1, $\infty$) \\
real & $\mathbb{R}$ & any number in (-$\infty$, $\infty$)\\
float & $F$ & floating point numbers within (approx.) (-1.798$*$10$^{308}$, 1.798$*$10$^{308}$) \\
object & $X$ & an abstract representation of a collection of the above, similar to tuple \\
array & $a$ & a collection of a single data type, $n$ entries of the form ($a_0$, $a_1$, $a_2$, ... ,$a_{n-1}$) \\
list & $L$ & an array of non-fixed length \\
\bottomrule
\end{tabular} 
\end{center}

\noindent
This specification of PhysioCompanion uses some derived data types: sequences, strings, and
tuples. Sequences are lists filled with elements of the same data type. Strings
are sequences of characters. Tuples contain a list of values, potentially of
different types. In addition, PhysioCompanion uses functions, which
are defined by the data types of their inputs and outputs. Local functions are
described by giving their type signature followed by their specification.

\section{Module Decomposition}

%Cieran check plssss
The following table is taken directly from the Module Guide document for this project.

\begin{table}[h!]
\centering
\begin{tabular}{p{0.3\textwidth} p{0.6\textwidth}}
\toprule
\textbf{Level 1} & \textbf{Level 2}\\
\midrule

{Hardware-Hiding Module} & ~ \\
\midrule

\multirow{4}{0.3\textwidth}{Behaviour-Hiding Module}& User Interface Module\\
& Home Module\\
& Main Module\\
& Draw Module \\
\midrule

\multirow{5}{0.3\textwidth}{Software Decision Module} & Database Management Module\\
& Generator Module\\
& Metrics Module\\
& Analyze Module\\
& Feedback Module\\
\bottomrule

\end{tabular}
\caption{Module Hierarchy}
\label{TblMH}
\end{table}


\section{MIS}
This section describes individual module specifications.\\

\subsection{Module - Account Creation}
This module enables the creation of a new account on the application. It uses a database to store any information that 
is essential for other modules to perform their actions. A database was used for its ease of manipulation and its 
comprehensive, multi-functional nature. 

\noindent The database tables related to this module are:
\begin{enumerate}
  \item User data (\texttt{users}): This table consists of information such as name,
  role, email, birthday, license\_no/invite\_code, createdAt and physioID.
  \item Invite\_codes (\texttt{inviteCodes}): This table contains all valid invite codes that will be used to validate 
  patient account creation on the application.
  \item Valid licenses (\texttt{validLicenses}): This is used for verification and validation purposes for a physiotherapist
  account creation.
\end{enumerate}  


\subsubsection{Uses}

\noindent \underline{\textbf{\textit{authService.ts - Variables}}}
\begin{description}[style=nextline, leftmargin=2.5cm]
\item[\texttt{auth}] 
    Firebase Authentication instance imported from "./firebase". Used for all authentication operations including email checking, sign-in, sign-out, and account creation.

\item[\texttt{db}] 
    Firestore database instance imported from "./firebase". Used for reading and writing user data, license validation, and invite code management.
\end{description}

\noindent \underline{\textbf{\textit{profileActivity.ts - Variables}}}\\
None

\noindent \underline{\textbf{\textit{authService.ts - Types}}}\\
None\\
\noindent \underline{\textbf{\textit{profileActivity.ts - Types}}}\\
None\\


\noindent \underline{\textbf{\textit{authService.ts - Functions}}}
\begin{description}[style=nextline, leftmargin=2.5cm]
\item[\texttt{createUserWithEmailAndPassword}] 
    Firebase Auth function from "firebase/auth". Used to create new user accounts in Firebase Authentication for both physiotherapist and patient registration.

\item[\texttt{doc}] 
    Firestore function from "firebase/firestore". Used to create document references for users, validLicenses, and inviteCodes collections.

\item[\texttt{getDoc}] 
    Firestore function from "firebase/firestore". Used to fetch user documents, license documents, and invite code documents from Firestore.

\item[\texttt{setDoc}] 
    Firestore function from "firebase/firestore". Used to create new user documents in Firestore during registration.

\item[\texttt{updateDoc}] 
    Firestore function from "firebase/firestore". Used to mark invite codes as used after patient registration.

\item[\texttt{serverTimestamp}] 
    Firestore function from "firebase/firestore". Used to set createdAt timestamps on new user documents.
\end{description}

\noindent \underline{\textbf{\textit{profileActivity.ts - Functions}}}
\begin{description}
\item[\texttt{getPhysioInviteCode(physioID)}] 
    Retrieves the invite code associated with a physiotherapist from the inviteCodes collection. Used during patient account creation to link new accounts to the correct physiotherapist.
\end{description}

\subsubsection{Syntax}

\textbf{Exported Constants}\\
\noindent \textbf{\textit{\underline{authService.ts}}}
\begin{description}[style=nextline, leftmargin=2.5cm]

\item[\texttt{export const registerPatient}] 
    Async function constant that registers a new patient account. Validates invite code, checks email uniqueness, creates Firebase Auth user, and stores patient data in Firestore.

\item[\texttt{export const registerPhysio}] 
    Async function constant that registers a new physiotherapist account. Validates license number format and existence, checks email uniqueness, creates Firebase Auth user, and stores physio data in Firestore.

\item[\texttt{export const checkEmailExists}] 
    Async function constant that checks if an email address is already registered in the system by querying the Firestore users collection.

\item[\texttt{export const licenseFormatValid}] 
    Async function that validates the format of a physiotherapist license number against the pattern XX-123456.

\end{description}

\noindent \underline{\textbf{\textit{profileActivity.ts}}} \\
None\\


\noindent \textbf{Exported Access Programs} 

\noindent \underline{\textbf{\textit{authService.ts}}} \\

\begin{longtable}{
|>{\columncolor{softgray}}p{4.2cm}|
>{\columncolor{white}}p{4.2cm}|
>{\columncolor{white}}p{4.2cm}|
}
\hline
\rowcolor{headerblue}
\textcolor{white}{\textbf{Input}} & 
\textcolor{white}{\textbf{Output}} & 
\textcolor{white}{\textbf{Exceptions}} \\
\hline
\endfirsthead

\hline
\rowcolor{headerblue}
\textcolor{white}{\textbf{Input}} & 
\textcolor{white}{\textbf{Output}} & 
\textcolor{white}{\textbf{Exceptions}} \\
\hline
\endhead

\multicolumn{3}{c}{\textit{Continued on next page}} \\
\endfoot
\endlastfoot

% checkEmailExists
\multicolumn{3}{|>{\columncolor{lightblue}}p{13.5cm}|}{\textbf{checkEmailExists(email):} Checks if an email is already registered in Firebase Authentication.} \\
\hline
\rowcolor{white}
\texttt{email: string} - User's email address &
\texttt{Promise<boolean>} - Returns true if email already exists, false otherwise &
\texttt{Error} - Throws error if there is an error checking email. \\
\hline

% licenseFormatValid
\multicolumn{3}{|>{\columncolor{lightblue}}p{13.5cm}|}{\textbf{licenseFormatValid(licenseNumber):} Validates physiotherapist license number format (XX-123456).} \\
\hline
\rowcolor{white}
\texttt{licenseNumber: string} - PT license number &
\texttt{boolean} - Returns true if format is valid, false otherwise &
\texttt{Error} - Throws error if there is an error checking licenseNumber. \\
\hline

% registerPhysio
\multicolumn{3}{|>{\columncolor{lightblue}}p{13.5cm}|}{\textbf{registerPhysio:} Registers a new physiotherapist account into the databese after license validation} \\
\hline
\rowcolor{white}
\texttt{params: \{name: string, email: string, password: string, licenseNumber: string\}} &
\texttt{Promise<\{uid: string, role: "physio"\}>} - Returns user ID and role &
\texttt{Error} - Invalid license format, license not verified, Firebase Auth error \\
\hline

% registerPatient
\multicolumn{3}{|>{\columncolor{lightblue}}p{13.5cm}|}{\textbf{registerPatient:} Registers a new patient account into the database.} \\
\hline
\rowcolor{white}
\texttt{params: \{name: string, email: string, password: string, inviteCode: string, birthday: string\}} &
\texttt{Promise<\{uid: string, role: "patient"\}>} - Returns user ID and role, marks invite code as used &
\texttt{Error} - Invalid invite code, inactive/used code, missing physio link, or Firebase Auth error \\
\hline
\caption{Account Creation Exported functions} \label{sec: accountexpprogs}\\
\end{longtable}

\noindent \underline{\textbf{\textit{profileActivity.ts}}} \\
None\\


\noindent \textbf{Exported Access Programs} \hypertarget{acexpprogs}\\
\texttt{p}: (name, email, password, licenseNumber, birthday)\\
\texttt{pt}: (name, email, password, inviteCode, birthday)\\
\texttt{user}: p or pt \\
\texttt{e}: user.email \\
\texttt{n}: user.name \\
\texttt{pw}: user.password\\
\texttt{lno}: user.licenseNumber\\
\texttt{code}: user.inviteCode \\
\texttt{bday}: user.birthday\\
\texttt{cat}: user.createdAt \\
\texttt{pid}: user.physioId (for \texttt{pt} only)\\
\texttt{v}: user.verified (for \texttt{p} only) \\
\texttt{u}: user.uid \\
\texttt{del}: user.deleted (for \texttt{pt} only)\\
\textit{db}: \texttt{users} table in database \\
\textit{db'}: newer version of \texttt{users} table \\
$i, j, k, l, m, n \in \{0-9\}$ \\
$a, b \in \{A-Z\}$ \\
\newline
\textit{E(e)}: 'checkEmailExists(e)' \\
\textit{L(l)}: 'licenseFormatValid(l)' \\
\textit{Phy(p)}: 'registerPhysio(p)'\\
\textit{Pt(pt)}: 'registerPatient(pt)'\\


\noindent \underline{Function definitions:} \label{sec: ACfuncdef} \\
\noindent $E(e) \implies e \in db.email$ \\
\newline
\noindent $L(l) \implies $ 
$\exists$ $l$ $|$ $l = (a + b $ $+$ $'-' $ + 'i' + 'j' + 'k' + 'l' + 'm' + 'n') $\in validLicenses$ \\
\newline
$Phy(p) \implies db' = db \cup \texttt{users}\mapsto\{{uid \mapsto \{cat, e, lno, n, role:'physio', v\}}\} $\\
\newline
$Pt(pt) \implies db' = db \cup \texttt{users}\mapsto\{{uid \mapsto \{cat, del, e, code, n, pid, role:'patient'\}}\} $\\


\subsubsection{Semantics}
\noindent \textbf{State Variables} \\
See \hyperlink{acexpprogs}{here} for a list of additional state variables.\\
\texttt{first}: user 1\\
\texttt{second}: user 2\\
\texttt{first\_uid}: uid of user 1\\
\texttt{second\_uid}: uid of user 2\\

\noindent \textbf{State Invariant} \\
$\forall$ $first, second \in \texttt{db}, first\_uid \cap second\_uid = \emptyset $\\
All users have different uids.\\

\subsubsection{Environment Variables}
\textit{screen}: The application's display to the user \\
\textit{touchscreen}: Users will click on the screen to select their desired action

\subsubsection{Assumptions}
\begin{itemize}
  \item Patients with the same name have different birthdays.
  \item Only one person can access user information in the database at a time since many values that are exported are asynchronous values.
\end{itemize}

\subsubsection{Access Routine Semantics - Exported}
\textbf{Exported access routines} \label{sec:acARS}\\
See \hyperlink{acexpprogs}{here} for state variables being used.
\newline
\textit{\underline{checkEmailExists: E(e)}}\\
\textbf{Pre-condition:}
  e $\in db.email$\\
\textbf{Function:}
  E(e) = true \\
\textbf{Post-condition:}
  Returns true if email already exists in database(\texttt{users}), false otherwise.\\
\textit{\underline{licenseFormatValid: L(l)}}\\
\textbf{Pre-condition:}
  l $\in db.validLicenses$\\
\textbf{Function:}
  L(l) = true \\
\textbf{Post-condition:}
  Returns true if license exists in database(\texttt{validLicenses}), false otherwise.\\
\textit{\underline{registerPhysio: Phy(p)}}\\
\textbf{Pre-condition:}
  $(p.name \cap p.email \cap p.pw \cap p.lno \cap p.bday \neq \emptyset)\cap (L(p.lno)) \cap (\neg (E(p.email)))$\\
\textbf{Function:}
   Phy(p)\\
\textbf{Post-condition:}
  $p \in \texttt{db}$\\
\textit{\underline{registerPatient: Pt(pt)}}\\
\textbf{Pre-condition:}
  $(p.name \cap p.email \cap p.pw \cap p.code \cap p.bday \neq \emptyset) \cap (\neg (E(p.email))) $\\
\textbf{Function:}
   Pt(pt)\\
\textbf{Post-condition:}
  $p \in \texttt{db}$\\

\subsubsection{Local Functions}

\begin{longtable}{
|>{\columncolor{softgray}}p{4.2cm}|
>{\columncolor{white}}p{4.2cm}|
>{\columncolor{white}}p{4.2cm}|
}
\hline
\rowcolor{headerblue}
\textcolor{white}{\textbf{Input}} & 
\textcolor{white}{\textbf{Output}} & 
\textcolor{white}{\textbf{Exceptions}} \\
\hline
\endfirsthead

\hline
\rowcolor{headerblue}
\textcolor{white}{\textbf{Input}} & 
\textcolor{white}{\textbf{Output}} & 
\textcolor{white}{\textbf{Exceptions}} \\
\hline
\endhead

\multicolumn{3}{c}{\textit{Continued on next page}} \\
\endfoot
\endlastfoot

\multicolumn{3}{|>{\columncolor{lightblue}}p{13.5cm}|}{\textbf{passwordCheck():} Internal helper function for password validation.} \\
\hline
\rowcolor{white}
\texttt{params: \{password: string\}} &
\texttt{string | null} &
\begin{tabular}{@{}l@{}}
• Returns "Password\\ must be at least \\6 characters long"\\ if length $<$ 6\\
• Returns "Password\\ must contain at least\\ one number" if no\\ numbers present\\
• Returns null if valid
\end{tabular} \\
\hline

\multicolumn{3}{|>{\columncolor{lightblue}}p{13.5cm}|}{\textbf{validateBirthday():} Internal helper function for birthday format validation.} \\
\hline
\rowcolor{white}
\texttt{params: \{birthday: string\}} &
\texttt{string | null} &
\begin{tabular}{@{}l@{}}
• Returns "Please enter \\birthday in YYYY\\-MM-DD format" if \\format entered is \\invalid\\
• Returns null if valid
\end{tabular} \\
\hline

\caption{Account creation module: Local functions}
\end{longtable}
\noindent \textit{Note: CreatePatient and CreatePhysio are in the mobile/app/(auth) folder. CreatePatient and registerPatient have an important distinction although they seem similar. 
CreatePatient is the React component which manages the interface, while registerPatient handles the backend logic for
making a new patient account. CreatePhysio and registerPhysio have the same distinction.} \\

\noindent \textbf{Local Routines (definitions)}\\
\texttt{p}: (name, email, password, licenseNumber, birthday)\\
\texttt{pt}: (name, email, password, inviteCode, birthday)\\
\texttt{user}: p or pt\\
\texttt{e}: user.email \\
\texttt{n}: user.name \\
\texttt{pw}: user.password\\
\texttt{lno}: user.licenseNumber\\
\texttt{code}: user.inviteCode\\
\texttt{bday}: user.birthday\\
\textit{Pc(pw)}: 'passwordCheck(pw)'\\
\textit{Vb(bday)}: 'validateBirthday(bday)' \\
\texttt{null}: null value\\


\noindent \underline{Function definitions}: \\
$Pc(pw) \implies pw.length \geq 6 \cap \exists$ $d$ $|$ $d \in \{0 - 9\} \cap d \in pw$ \\
\texttt{Date} = 'YYYY-MM-DD' $|$ 'YYYY' $\in$ \{0000 - 9999\} $\cap$ 'MM' $\in$ \{01 - 12\} $\cap$ 'DD' $\in$ \{01 - 31\} \\
$Vb(bday) \implies bday \in Date $ \\
\textit{Note: See \ref{sec:ui} for CreatePhysio and CreatePatient interface specifications.}

\subsubsection{Access Routine Semantics - Local}
\underline{\textit{Vb(bday)}}:\\
\textbf{Pre-condition}: \texttt{bday} $\notin \emptyset$\\
\textbf{Function}: Vb(bday)\\
\textbf{Post-condition}: returns \texttt{null} if valid \\
\underline{\textit{Pc(pw)}}:\\
\textbf{Pre-condition}: pw $\notin \emptyset$\\
\textbf{Function}: Pc(pw)\\
\textbf{Post-condition}: returns \texttt{null} if valid \\ 
%%%%%%%%%%%%%%%%%%%%%%%%%%%%%%%%%%%%%%%%%%%%%%%%%%%%%%%%%%%%%%%%%%%%
\subsection{Module - Valid user}
This module's purpose is to check whether a potential user is authorized to create an account.
For feasible implementation purposes, a file with hard coded invite codes and license numbers will be used.
The actions accomplishing the module's purpose are within the registerPhysio and registerPatient functions,
which contain internal validation of inviteCode and licenseNumber.
\subsubsection{Uses}
\noindent \underline{\textbf{\textit{authService.ts}}} \\
\begin{description}[style=nextline, leftmargin=4cm]

\item[\texttt{params}] 
    Input object containing patient registration data. Type: \texttt{\{name: string, email: string, password: string, inviteCode: string, birthday: string\}}.

\item[\texttt{code}] 
    Normalized invite code. Type: string. Computed as \texttt{params.inviteCode.trim().\\toUpperCase()}.

\item[\texttt{inviteRef}] 
    Firestore document reference to \texttt{inviteCodes/code}. Type: \texttt{DocumentReference}.

\item[\texttt{inviteSnap}] 
    Firestore document snapshot. Type: \texttt{DocumentSnapshot}.

\item[\texttt{invite}] 
    Data from invite code document. Contains fields: \texttt{active} (boolean), \texttt{used} (boolean), \texttt{physioId} (string).

\item[\texttt{physioId}] 
    \texttt{uid} of the physiotherapist that provided invite code. Type: string.

\item[\texttt{cred}] 
    Firebase Auth user credential. Type: \texttt{UserCredential}.

\item[\texttt{uid}] 
    Firebase Authentication unique user ID. Type: string.

\item[\texttt{db}] 
    Firestore database instance. Imported from \texttt{./firebase}.

\item[\texttt{auth}] 
    Firebase Authentication instance. Imported from \texttt{./firebase}.

\item[\texttt{licenseNumber}] 
    State variable managing license number input. Type: string. Initialized as empty string.

\item[\texttt{lic}] 
    Normalized license number. Type: string. Computed as \texttt{params.licenseNumber.\\trim().toUpperCase()}.

\item[\texttt{licSnap}] 
    Firestore document snapshot from \texttt{validLicenses} collection. Type: \texttt{DocumentSnapshot}.

\end{description}
\subsubsection{Syntax}
\textbf{Exported Constants}\\
See \textbf{\textit{authService.ts}} exported constants in \ref{sec:validexp}. \\
For interface elements, see \ref{sec:ui}.\\
\newline
\noindent \textbf{Exported Access Programs} \\
See Table \ref{sec: accountexpprogs} for \textit{\textbf{registerPatient}} and \textit{\textbf{registerPhysio}}
details respectively.\\
\newline
\underline{\textit{Function definitions:}}\\
This module has components that have been defined in the Account Creation module. See \hyperref[sec: ACfuncdef]{\textit{function definitions in \ref*{sec: ACfuncdef}}}.

\subsubsection{Semantics}
\textbf{State Variables}\\
See \ref{sec:ui} for UI state variables. \\
\texttt{pt}: a patient user type\\
\texttt{p}: a physio user type \\
\texttt{code}: a patient invite code \\
\texttt{lno}: a PT's license number \\
\texttt{users}: \texttt{users} table in database\\

\noindent \textbf{State Invariant}\\
$ code['used'] = true \implies \exists$ $pt |$ $pt.inviteCode == code$\\
$ \exists p.uid \in \texttt{users} \implies p.licenseNumber \in \texttt{validLicenses}$ 
\subsubsection{Environment state variables}
\textit{Not applicable}
\subsubsection{Assumptions}
\begin{itemize}
    \item All practicing physiotherapists have an active and valid license.
    \item All active invite codes are used by new and active patients only.
    \item Invite codes and license numbers are unique.
\end{itemize}

\subsubsection{Access Routine Semantics - Exported}
\textbf{Exported access routines} \\
See \hyperref[sec:acARS]{account creation semantics} for details of respective functions.

\subsubsection{Local Routines}
Portions that are pertinent to the module (from \underline{\textbf{\textit{authService.ts}}} \\) are elaborated on below. \\

\noindent \textbf{registerPatient: Internal Validation Steps}\\
The function performs the following sequential validations on the invite code:

\begin{enumerate}
\item \textbf{Existence Check}: Query Firestore \texttt{inviteCodes} collection for document with ID = \texttt{normalize(inviteCode)}. If not found, throw "Invalid invite code."

\item \textbf{Active Check}: If \texttt{invite.active $\neq$ true}, throw "Invite code inactive."

\item \textbf{Usage Check}: If \texttt{invite.used = true}, throw "Invite code already used."

\item \textbf{Physio Link Check}: If \texttt{physioId = null} or does not exist, throw "Invite code missing physio link."

\end{enumerate}

Only after all validations pass does the function proceed to:
\begin{enumerate}
\item Create Firebase Auth user
\item Create Firestore user document with \texttt{role:"patient"} and \texttt{uid}
\item Mark invite code as used with \texttt{usedBy:uid}
\end{enumerate}

\noindent \textbf{registerPhysio: Internal Validation Steps}\\
The function performs the following sequential validations on the license number:

\begin{enumerate}
\item \textbf{Format Check}: Validate license number format using \texttt{licenseFormatValid()}. License must match pattern: two uppercase letters, a dash, and six digits (e.g., ON-123456). If format is invalid, throw "Invalid license format (example: ON-123456)."

\item \textbf{Existence Check}: Query Firestore \texttt{validLicenses} collection for document with ID equal to the normalized license number. If not found, throw "License not verified."

\item \textbf{Active Check}: If the license document has \texttt{active} field not equal to true, throw "License not verified."

\end{enumerate}

Only after all validations pass does the function proceed to:
\begin{enumerate}
\item Create Firebase Auth user using \texttt{createUserWithEmailAndPassword()}
\item Create Firestore user document with \texttt{role:"physio"}, \texttt{verified:true}, and \texttt{uid}
\item Return \texttt{\{uid, role:"physio"\}}
\end{enumerate}

\subsubsection{Access Routine Semantics - Local}
NA\\
See \href{sec:acARS}{Account creation access routine semantics}



%%%%%%%%%%%%%%%%%%%%%%%%%%%%%%%%%%%%%%%%%%%%%%%%%%%%%%%%%%%%%%%%%%%%%%%
\subsection{Module - User Account} \label{section:user account}

This module's purpose is facilitate data manipulation within the user database tables. It was created as a separate component
to simplify implementation.
%%%%%%%%%%%%%%%%%%%%%%%%%%%\\\\\\\\\
Account creation generates a single instance of one of the following objects:\\
\begin{itemize}
    \item \textbf{Patient object:} 
        \texttt{\{uid, email, name, role: "patient", birthday?, physioId?, inviteCode?, createdAt?, deleted?\}}
    
    \item \textbf{Physio object:}
        \texttt{\{uid, email, name, role: "physio", licenseNumber?, verified?, createdAt?\}}
\end{itemize}

\subsubsection{Uses}
\underline{\textbf{\textit{useraccount.ts}}}
\begin{description}
\item[\texttt{uid: string;}] 
    Unique identifier for the user. Used for fetching user data, updating user information, deleting user accounts, and associating patients with physiotherapists.

\item[\texttt{email: string;}] 
    User's email address for authentication and communication. Used for email existence verification, credential validation, and user authentication.

\item[\texttt{name: string;}] 
    User's full name. Used for searching users by name, retrieving user database information by name, and patient account deletion verification.

\item[\texttt{role: "patient" | "physio";}] 
    User's role in the system. Used for filtering patients by physiotherapist and categorizing users during database information retrieval.

\item[\texttt{physioId?: string;}] 
    Reference to the physiotherapist assigned to a patient. Used for retrieving all patients belonging to a specific physiotherapist.

\item[\texttt{deleted?: boolean;}]
    Soft delete flag indicating whether the user account has been deleted. Used as a filter when querying active patients and during user account deletion operations.

\item[\texttt{birthday?: string;}] 
    User's birthday. Used as part of patient profile information.

\item[\texttt{licenseNumber?: string;}] 
    Professional license number for physiotherapists. Used as part of physiotherapist profile information.

\item[\texttt{verified?: boolean;}] 
    Flag indicating whether a physiotherapist's license has been verified. Used in physiotherapist profile information.

\item[\texttt{inviteCode?: string;}] 
    Invitation code for patient registration. Used as profile information and account creation validation.

\item[\texttt{createdAt?: any;}]
    Timestamp of user account creation. Defined in the interface but not set or read by current methods. Recommended for future audit trail implementation.

\end{description}
%%%%%%%%%%%%\\\\\\\\\\\\\\\\

\subsubsection{Syntax}

\textbf{Exported Constants}\\
None\\

\noindent \textbf{Exported Access Programs}\\
\noindent \underline{\textbf{\textit{useraccount.ts}}}

\begin{longtable}{
|>{\columncolor{softgray}}p{4.2cm}|
>{\columncolor{white}}p{4.2cm}|
>{\columncolor{white}}p{4.2cm}|
}
\hline
\rowcolor{headerblue}
\textcolor{white}{\textbf{Input}} & 
\textcolor{white}{\textbf{Output}} & 
\textcolor{white}{\textbf{Exceptions}} \\
\hline
\endfirsthead

\hline
\rowcolor{headerblue}
\textcolor{white}{\textbf{Input}} & 
\textcolor{white}{\textbf{Output}} & 
\textcolor{white}{\textbf{Exceptions}} \\
\hline
\endhead

\multicolumn{3}{c}{\textit{Continued on next page}} \\
\endfoot
\endlastfoot

% validateCredentials
\multicolumn{3}{|>{\columncolor{lightblue}}p{13.5cm}|}{\textbf{validateCredentials(email, password):} Authenticates user with Firebase Auth.} \\
\hline
\rowcolor{white}
\texttt{email: string, password: string} &
\texttt{Promise<boolean>} &
\texttt{None - returns false on failure} \\
\hline

% getUserData
\multicolumn{3}{|>{\columncolor{lightblue}}p{13.5cm}|}{\textbf{getUserData(uid):} Retrieves user data from Firestore by UID.} \\
\hline
\rowcolor{white}
\texttt{uid: string} &
\texttt{Promise<UserData | null>} &
\texttt{UserAccountField} \texttt{EmptyError} \\
\hline

% getUserByEmail
\multicolumn{3}{|>{\columncolor{lightblue}}p{13.5cm}|}{\textbf{getUserByEmail(email):} Retrieves user data from Firestore by email.} \\
\hline
\rowcolor{white}
\texttt{email: string} &
\texttt{Promise<UserData | null>} &
\texttt{None} \\
\hline

% emailExists
\multicolumn{3}{|>{\columncolor{lightblue}}p{13.5cm}|}{\textbf{emailExists(email):} Checks if an email is already registered.} \\
\hline
\rowcolor{white}
\texttt{email: string} &
\texttt{Promise<boolean>} &
\texttt{None} \\
\hline

% updateUser
\multicolumn{3}{|>{\columncolor{lightblue}}p{13.5cm}|}{\textbf{updateUser(uid, updates):} Updates user data in Firestore.} \\
\hline
\rowcolor{white}
\texttt{uid: string, updates: Partial<UserData>} &
\texttt{Promise<boolean>} &
\texttt{UserAccountNot} \texttt{FoundError} \\
\hline

% deleteUser
\multicolumn{3}{|>{\columncolor{lightblue}}p{13.5cm}|}{\textbf{deleteUser(uid):} Soft-deletes a user by setting deleted flag to true.} \\
\hline
\rowcolor{white}
\texttt{uid: string} &
\texttt{Promise<boolean>} &
\texttt{None} \\
\hline

% deleteUserByEmail
\multicolumn{3}{|>{\columncolor{lightblue}}p{13.5cm}|}{\textbf{deleteUserByEmail(email):} Soft-deletes a user by email address.} \\
\hline
\rowcolor{white}
\texttt{email: string} &
\texttt{Promise<boolean>} &
\texttt{UserAccountNot} \texttt{FoundError} \\
\hline

% getPatientsByPhysio
\multicolumn{3}{|>{\columncolor{lightblue}}p{13.5cm}|}{\textbf{getPatientsByPhysio(physioId):} Retrieves all active patients assigned to a physiotherapist.} \\
\hline
\rowcolor{white}
\texttt{physioId: string} &
\texttt{Promise<UserData[]>} &
\texttt{None} \\
\hline

% getUsersByName
\multicolumn{3}{|>{\columncolor{lightblue}}p{13.5cm}|}{\textbf{getUsersByName(name):} Searches users by partial name match.} \\
\hline
\rowcolor{white}
\texttt{name: string} &
\texttt{Promise<UserData[]>} &
\texttt{None} \\
\hline

% getAllUsers
\multicolumn{3}{|>{\columncolor{lightblue}}p{13.5cm}|}{\textbf{getAllUsers():} Retrieves all users from Firestore.} \\
\hline
\rowcolor{white}
\texttt{None} &
\texttt{Promise<UserData[]>} &
\texttt{None} \\
\hline

% getUserdbInfo
\multicolumn{3}{|>{\columncolor{lightblue}}p{13.5cm}|}{\textbf{getUserdbInfo(name, birthday?):} Retrieves users by name and optional birthday, separated by role.} \\
\hline
\rowcolor{white}
\texttt{name: string, birthday?: string} &
\texttt{Promise<\{patients: UserData[], physios: UserData[]\}>} &
\texttt{UserAccountNot} \texttt{FoundError} \\
\hline

% PTaccountDelete
\multicolumn{3}{|>{\columncolor{lightblue}}p{13.5cm}|}{\textbf{PTaccountDelete(physioEmail, patientName, patientEmail):} Allows physiotherapist to soft-delete a linked patient account.} \\
\hline
\rowcolor{white}
\texttt{physioEmail: string, patientName: string, patientEmail: string} &
\texttt{Promise<void>} &
\begin{tabular}{@{}l@{}}
\texttt{UserAccountField}\\ \texttt{EmptyError} \\
\texttt{UserAccountNot} \\ \texttt{FoundError} \\
\end{tabular} \\
\hline

% usernamePwMatch
\multicolumn{3}{|>{\columncolor{lightblue}}p{13.5cm}|}{\textbf{usernamePwMatch(email, password):} Verifies email and password combination.} \\
\hline
\rowcolor{white}
\texttt{email: string, password: string} &
\texttt{Promise<boolean>} &
\texttt{UserAccountLogin} \texttt{MatchError} \\
\hline

% authenticateUser
\multicolumn{3}{|>{\columncolor{lightblue}}p{13.5cm}|}{\textbf{authenticateUser(email, password):} Authenticates user and returns user data if successful.} \\
\hline
\rowcolor{white}
\texttt{email: string, password: string} &
\texttt{Promise<UserData | null>} &
\texttt{None} \\
\hline

% initializeSystem
\multicolumn{3}{|>{\columncolor{lightblue}}p{13.5cm}|}{\textbf{initializeSystem():} Initializes the system (placeholder for future setup).} \\
\hline
\rowcolor{white}
\texttt{None} &
\texttt{Promise<void>} &
\texttt{None} \\
\hline

\caption{UserAccount module: Exported access programs}
\end{longtable}

\subsubsection{Semantics}
\noindent \textbf{State Variables} \\
\underline{\textbf{\textit{useraccount.ts}}}
\begin{description}
\item[\texttt{private usersCollection: string}] 
    Internal state variable storing the Firestore collection name for user documents.

\item[\texttt{private validLicensesCollection: string}] 
    Internal state variable storing the Firestore collection name for physiotherapist licenses.

\item[\texttt{private inviteCodesCollection: string}] 
    Internal state variable storing the Firestore collection name for invitation codes.
\end{description}

\noindent \textbf{State Invariant} \\
\texttt{users}: 'users' database table\\
\texttt{u1}: user 1\\
\texttt{u2}: user 2\\
\texttt{uid}: uid\\
\texttt{email}: email\\
\texttt{Phy}: set of all users that have role = 'physio' \\
\texttt{Pt}: set of all users that have role = 'patient' \\
\textit{U}: set of all UserData objects\\
\textit{E}: set of valid emails
\texttt{updates}: part of UserData object 


\noindent $\forall$ $u1, u2 \in users, u1 \neq u2 \implies u1.email \neq u2.email$ \\
$\forall$ $u1$ $\in$ $Phy, u1.physioId = \emptyset$ \\
$\forall$ $u1$ $\in$ $Phy, u1.verified = true$ \\
%$\forall$ $u1, u2$ $\in users$ $|$ $u1 \in Pt \cap u2 \in Phy,$ $u1.physioId = u2.uid$ \\
$\forall$ $u1 \in Pt,$ $u1.deleted = true \implies u1.uid$ $\in \emptyset$ \\
$\forall$ $u1 \in P, u2 \in Phy,$ $\exists$ $u1.physioId = u2.uid$ \\
$\forall$ $id \in \{users.uid\}, $$\exists$ $u1| u1.uid = id \cap u1 \in users$ \\
$\forall$ $u1 \in users \cap u1 \in Phy \implies u1 \notin P$ \\
$\forall$ $u2 \in users \cap u2 \in Pt \implies u2 \notin Phy$ \\


\subsubsection{Environment Variables}
Not applicable

\subsubsection{Assumptions}
\begin{itemize}
  \item If two people have the same name, then it is assumed they have different different roles.
  \item No two PTs have the same license numbers.
\end{itemize}

\subsubsection{Access Routine Semantics - Exported}
\textit{Note: Auth functions access routine semantics are not outlined here.}\\
\underline{\textit{getUserData(uid)}}\\
\textbf{Pre-condition}: \texttt{uid} $\notin \emptyset$\\
\textbf{Function}: getUserData(uid)\\
\textbf{Post-condition}: returns UserData object | \texttt{uid} = user.uid\\

\noindent \underline{\textit{getUserByEmail(email)}}\\
\textbf{Pre-condition}: \texttt{email} $\notin \emptyset$\\
\textbf{Function}: getUserByEmail(email)\\
\textbf{Post-condition}: returns UserData object | \texttt{email} = user.email\\

\noindent\underline{\textit{emailExists(email)}}\\
\textbf{Pre-condition}: email is a valid email\\
\textbf{Function}: emailExists(email)\\
\textbf{Post-condition}: $True \implies email \in E$\\

\noindent\underline{\textit{updateUser(uid, updates)}}\\
\textbf{Pre-condition}: $user \in users | \exists user.uid = \texttt{uid}$\\
\textbf{Function}: emailExists(email)\\
\textbf{Post-condition}: $True \implies email \in E$\\

\noindent\underline{\textit{deleteUser(uid)}}\\
\textbf{Pre-condition}: $user \in users$ \\
\textbf{Function}: deleteUser(uid)\\
\textbf{Post-condition}: $users' = users - (user[uid])$\\

\noindent\underline{\textit{deleteUserByEmail(email)}}\\
\textbf{Pre-condition}: $user \in users$ \\
\textbf{Function}: deleteUserByEmail(email)\\
\textbf{Post-condition}: $users' = users - (user[uid])$\\


\subsubsection{Local Functions}
Not applicable. All helper functions are defined as private class functions. 

\subsubsection{Access Routine Semantics - Local}
Not applicable

%%%%%%%%%%%%%%%%%%%%%%%%%%%%%%%%%%%%%%%%%%%%%%%%%%%%%%%%%%%%%%%%%%%%%%%
\subsection{Module - Exercise Data}

This module manages exercises that will be performed by the patient. It consists of a data layer (\textit{exerciseData.ts})
and a UI layer (\textit{exercises.tsx}). The data layer provides database operations using the \texttt{userExercises} and \texttt{exercises} table.
The UI layer fulfills the purpose of interfacing between the user and the module and is discussed in
\ref{sec:ui}. This structure was chosen to enable:
\begin{itemize}
    \item High cohesion: all constituents serve the same business purpose
    \item Low coupling: there is no direct communication present between the front end and database. They communicate through the backend.
    \item Easy maintenance: low coupling allows for maintenance of a single 'part' of the module without affecting the rest of the module. 
\end{itemize}


\subsubsection{Uses}
\noindent \underline{\textbf{\textit{exerciseData.ts - Functions}}}
\begin{description}
\item[\texttt{collection}] 
    Firestore function from "firebase/firestore". Used to reference a database collection for exercise operations.

\item[\texttt{deleteDoc}] 
    Firestore function from "firebase/firestore". Used to delete an exercise document from the userExercises collection.

\item[\texttt{doc}] 
    Firestore function from "firebase/firestore". Used to create a document reference when fetching a single exercise.

\item[\texttt{getDoc}] 
    Firestore function from "firebase/firestore". Used to fetch a single exercise document by ID.

\item[\texttt{getDocs}] 
    Firestore function from "firebase/firestore". Used to fetch multiple exercise documents.

\item[\texttt{query}] 
    Firestore function from "firebase/firestore". Used to create filtered queries on exercise collections.

\item[\texttt{where}] 
    Firestore function from "firebase/firestore". Used to specify filter conditions (userid, title) in exercise queries.

\item[\texttt{addDoc}] 
    Firestore function from "firebase/firestore". Used to add a new exercise document to the userExercises collection.

\item[\texttt{updateDoc}] 
    Firestore function from "firebase/firestore". Used to update exercise fields (description, reps, sets, title) in an existing document.
\end{description}

\noindent \underline{\textbf{\textit{exerciseData.ts - Variables}}}
\begin{description}
\item[\texttt{db}] 
    Firestore database instance from "./firebase". Used as the database reference for all Firestore operations related to this file.
\end{description}

\noindent \underline{\textbf{\textit{exerciseData.ts - Types}}}
\begin{description}
\item[\texttt{Exercise}] 
    Type definition defined locally within this file. Represents the exercise data structure with properties: id, title, subtitle, description, demoText, enabled, tags, sets, reps.
\end{description}

\subsubsection{Syntax}

\noindent \textbf{Exported Constants}\\
None \\

\noindent \textbf{Exported Access Programs}\\
\textit{exerciseData.ts}
\begin{longtable}{
|>{\columncolor{softgray}}p{4.2cm}|
>{\columncolor{white}}p{4.2cm}|
>{\columncolor{white}}p{4.2cm}|
}
\hline
\rowcolor{headerblue}
\textcolor{white}{\textbf{Input}} & 
\textcolor{white}{\textbf{Output}} & 
\textcolor{white}{\textbf{Exceptions}} \\
\hline
\endfirsthead

\hline
\rowcolor{headerblue}
\textcolor{white}{\textbf{Input}} & 
\textcolor{white}{\textbf{Output}} & 
\textcolor{white}{\textbf{Exceptions}} \\
\hline
\endhead

\multicolumn{3}{c}{\textit{Continued on next page}} \\
\endfoot
\endlastfoot

% getExercisesById
\multicolumn{3}{|>{\columncolor{lightblue}}p{13.5cm}|}{\textbf{getExercisesById(uid):} Retrieves all exercises assigned to a specific user from the "userExercises" collection.} \\
\hline
\rowcolor{white}
\texttt{uid: string} - User ID of the patient &
\texttt{Promise<Exercise[] | undefined>} - Returns array of exercises or undefined if none found &
\texttt{Error} - Throws error if there is an issue fetching exercises from Firestore \\
\hline

% getExercise
\multicolumn{3}{|>{\columncolor{lightblue}}p{13.5cm}|}{\textbf{getExercise(exid):} Fetches a single exercise document by its ID.} \\
\hline
\rowcolor{white}
\texttt{exid: string} - Exercise document ID &
\texttt{Promise<Exercise | undefined>} - Returns exercise object or undefined if not found &
\texttt{Error} - Throws error if there is an issue fetching the exercise from Firestore \\
\hline

% getGeneralExercises
\multicolumn{3}{|>{\columncolor{lightblue}}p{13.5cm}|}{\textbf{getGeneralExercises():} Retrieves the master exercise library from the "exercises" collection.} \\
\hline
\rowcolor{white}
\texttt{None} &
\texttt{Promise<Exercise[] | undefined>} - Returns array of all general exercises or undefined if none found &
\texttt{Error} - Throws error if there is an issue fetching exercises from Firestore \\
\hline

% getSelectedExercises
\multicolumn{3}{|>{\columncolor{lightblue}}p{13.5cm}|}{\textbf{getSelectedExercises(uid):} Fetches the current list of exercises assigned to a specific user.} \\
\hline
\rowcolor{white}
\texttt{uid: string} - User ID of the patient &
\texttt{Promise<Exercise[] | undefined>} - Returns array of assigned exercises or undefined if none &
\texttt{Error} - Throws error if there is an issue fetching exercises from Firestore \\
\hline

% removeUserExercise
\multicolumn{3}{|>{\columncolor{lightblue}}p{13.5cm}|}{\textbf{removeUserExercise(uid, exerciseTitle):} Deletes an exercise assignment from a user's exercise list.} \\
\hline
\rowcolor{white}
\texttt{uid: string} - User ID of the patient, \texttt{exerciseTitle: string} - Title of the exercise to remove &
\texttt{Promise<void>} - Returns nothing on successful deletion &
\texttt{Error} - Throws error if there is an issue deleting the document from Firestore \\
\hline

% addUserExercise
\multicolumn{3}{|>{\columncolor{lightblue}}p{13.5cm}|}{\textbf{addUserExercise(uid, exercise):} Creates a new exercise assignment for a specific user.} \\
\hline
\rowcolor{white}
\texttt{uid: string} - User ID of the patient, \texttt{exercise: Exercise} - Exercise object to assign &
\texttt{Promise<void>} - Returns nothing on successful creation &
\texttt{Error} - Throws error if there is an issue adding the document to Firestore \\
\hline

% updateUserExercise
\multicolumn{3}{|>{\columncolor{lightblue}}p{13.5cm}|}{\textbf{updateUserExercise(uid, exerciseTitle, exerciseDescription, exerciseSets, exerciseReps):} Modifies an existing user exercise's instructions (description, sets, reps).} \\
\hline
\rowcolor{white}
\texttt{uid: string} - User ID of the patient, \texttt{exerciseTitle: string} - Title of the exercise, \texttt{exerciseDescription: string} - New description, \texttt{exerciseSets: string} - Number of sets (optional), \texttt{exerciseReps: string} - Number of reps (optional) &
\texttt{Promise<void>} - Returns nothing on successful update &
\texttt{Error} - Throws error if there is an issue updating the document in Firestore \\
\hline
\end{longtable}
\subsubsection{Semantics}

\noindent \textbf{State Variables} \hypertarget{svex} \\
\texttt{userexercises}: \texttt{userExercises} table in the database\\
\texttt{exercises}: \texttt{exercises} table in user database\\
\texttt{users}: \texttt{users} table in the database \\
\texttt{ue}: modified \texttt{exercise} object \\
\texttt{uue}: unmodifies \texttt{exercise} object\\
\texttt{assigned}: list of assigned exercises \\
\texttt{u1}: user \\
\texttt{exid}: exercise id \\
\texttt{ex}: exercise (object) \\
%\texttt{extitle}: exercise title\\
%\texttt{userid}: uid


\noindent \textbf{State Invariant} \\
$\forall$ $e1, e2 \in exercises, e1 \neq e2 \implies e1.title \neq e2. title \cap e1.description \neq e2.description$ \\
$\forall$ $e1 \in exercises,$ $\exists$ $e1.reps, e1.sets \in \mathbb{Z}^{+}$ \\
$\forall$ $e1, e2 \in userexercises, e1.title = e2.title \centernot \implies (e1.id = e2.id \cup 
e1.reps = e2.reps \cup e1.sets = e2.sets \cup e1.userid = e2.userid)$ \\
$\forall$ $ue, e1 \in userexercises$ $|$ $ue.userid = e1.userid \cap ue.title = e1.title, \exists ue.description \neq e1.description
\cup ue.sets \neq e1.sets \cup ue.reps \neq e1.reps$ \\
$\forall$ $user \in users, |assigned| \geq 0$ \\
\textit{\underline{Function definitions}}\\
\texttt{u1}: user\\
\texttt{uid}: uid (e.g. u1.uid)\\
\texttt{title}: exercise title \\
\textit{Ux}: 'getExercisesById(uid)'\\
\textit{Ge}: 'getExercise(exid)'\\
\textit{Gen}: 'getGeneralExercises()'\\
\textit{Sel}: 'getSelectedExercises(uid)'\\
\textit{R}: 'removeUserExercise(uid, ex.title)\\
\textit{A}: 'addUserExercise(uid, ex)'
\textit{Ux(uid)}: $\forall ex \in userexercises$ $Ux(uid) = [ex]$ $|$ $ex.userid = uid$\\
\textit{Ge(exid)}: $\forall$ $ex \in exercises, Ge(exid) = [ex]$ $|$ $ex.id = exid$\\ 
\textit{Gen()}: $Gen() \rightarrow \{ex\} \in exercises$ \\
\textit{Sel(uid)}: $\forall ex \in userexercises,$ $Sel(uid) = [ex]$ $|$ $ex.userid = uid$\\
\textit{R(uid, title)}: $userexercises' = userexercises - [ex]$ $|$ $ex.title = title \cap ex.userid = uid$\\
\textit{A(uid, ex)}: $userexercises' = userexercises \cup [ex]$ $|$ $ex \in exercises \cap (ex \in userexercises \cap ex.userid = uid) $ \\

\subsubsection{Environment Variables} 
Not applicable

\subsubsection{Assumptions}
\begin{itemize}
  \item For the same user, no two exercises with the same name are present. Conversely, no two 
  exercises with the same names can be assigned by a PT to one patient.
\end{itemize}

\subsubsection{Access Routine Semantics - Exported}
\texttt{db}: database\\
\texttt{userexercises}: \texttt{userExercises} table in the database\\
\texttt{exercises}: \texttt{exercises} table in user database\\
\texttt{users}: \texttt{users} table in the database \\
\texttt{ue}: modified \texttt{exercise} object \\
\texttt{assigned}: list of assigned exercises \\
\texttt{u1}: user \\
\texttt{exid}: exercise id \\
\texttt{ex}: exercise (object) \\
\newline
\textit{getExercisesById}\\
\textbf{Pre-condition}: $u1.uid = ex.userid \in \{userexercises.userid\}$ \\
\textbf{Function}: getExercisesById(u1.uid) \\
\textbf{Post-condition}: [ex]\\
\newline
\textit{getExercise}\\
\textbf{Pre-condition}: $exid \in \{exercises.id\} \cap exid \neq \emptyset$\\
\textbf{Function}: getExercises(ex.exid) \\
\textbf{Post-condition}: $ex \impliedby ex.id = exid $ \\
\newline
\textit{getGeneralExercises}\\
\textbf{Pre-condition}: $\exists exercises \in db $
\textbf{Function}: getGeneralExercises() \\
\textbf{Post-condition}: $[ex] \neq \emptyset$ \\
\newline
\textit{getSelectedExercises}\\
\textbf{Pre-condition}: $u1 \in users \cap u1.uid = ex.userid$ \\
\textbf{Function}: getSelectedExercises(u1.uid) \\
\textbf{Post-condition}: $[ex]$ \\
\newline
\textit{removeUserExercise}\\
\textbf{Pre-condition}: $u1 \in users \cap u1.uid = ex.userid \cap title = ex.title \in userexercises$
\textbf{Function}: getSelectedExercises(u1.uid, ex.title) \\
\textbf{Post-condition}: $user[exercises]' = user[exercises] - ex$ \\
\newline
\textit{addUserExercise}\\
\textbf{Pre-condition}: $u1 \in users \cap ex \in exercises$ \\
\textbf{Function}: addUserExercise(u1.uid, ex)\\
\textbf{Post-condition}: $user[exercises]' = user[exercises] + ex$\\
\newline
\textit{updateUserExercise}\\
\textbf{Pre-condition}: $u1 \in users \cap ex \in exercises \cap ex.userid = u1.uid$ \\
\textbf{Function}: addUserExercise(u1.uid, ex.title, ex.description, ex.sets, ex.reps)\\
\textbf{Post-condition}: $ue = ex$ $\cap$ $(ex(u1.uid, ue.title)$ $\cup$ $ex(u1.uid, ue.description)$
$\cup$ $ex(u1.uid, ue.sets)$ $\cup$ $ex(u1.uid, ue.reps))$

\subsubsection{Local Functions}
NA \\

\subsubsection{Access Routine Semantics - Local}
NA: they are UI related.


%%%%%%%%%%%%%%%%%%%%%%%%%%%%%%%%%%%%%%%%%%%%%%%%%%%%%%%%%%%%%%%%%%%%%%
\subsection{Module - User Activity}
\begin{comment}
This module handles video data of a patient's recorded exercise(s).

\subsubsection{Uses}
\begin{itemize}
  \item name
  \item exercise
  \item date
  \item duration
  \item recent\_video
  \item recent\_analysis
  \item overall\_analysis

\end{itemize}

\subsubsection{Syntax}

\noindent \textbf{Exported Constants} \\
\textit{analysis}: recent\_analysis\\
\textit{recent\_video}: recent\_video\\
\textit{analysis\_report}: analysis\_report\\

\noindent \textbf{Exported Access Programs} %TODO fix in and out
\begin{table}[h]
\centering
\rowcolors{2}{lavender}{white}
\setlength{\extrarowheight}{3pt}
\setlength{\arrayrulewidth}{1.2pt}

\begin{tabularx}{\textwidth}{|>{\raggedright\arraybackslash}X|>{\raggedright\arraybackslash}X|>{\raggedright\arraybackslash}X|>{\raggedright\arraybackslash}X|}
\hline
\rowcolor{darkpurple}
\color{white}\textbf{Access Routine Name} & \color{white}\textbf{Input} & \color{white}\textbf{Output} & \color{white}\textbf{Exceptions} \\
\hline

\rowcolor{lavender}
\multicolumn{4}{|p{\dimexpr\textwidth-2\tabcolsep}|}{\textbf{Description:} adds analysis results to user activity history} \\
\hline
\rowcolor{white}
add\_analysis\_to\newline\_user\_act & \makecell[lt]{\textit{user\_acc\_p}\\\textit{recent\_analysis}} & NA & \makecell[lt]{UserNotFoundError\\SaveError} \\
\hline
\rowcolor{lavender}
\multicolumn{4}{|p{\dimexpr\textwidth-2\tabcolsep}|}{\textbf{Description:} saves video recording for user account} \\
\hline
\rowcolor{white}
save\_video & \makecell[lt]{\textit{user\_acc\_p}\\\textit{recent\_video}} & NA & \makecell[lt]{UserNotFoundError\\UploadError\\VideoTooShortError} \\
\hline

\rowcolor{lavender}
\multicolumn{4}{|p{\dimexpr\textwidth-2\tabcolsep}|}{\textbf{Description:} retrieves analysis data for a specific date} \\
\hline
\rowcolor{white}
get\_analysis & \makecell[lt]{\textit{user\_acc\_p}\\\textit{date}} & \textit{recent\_analysis} & \makecell[lt]{UserNotFoundError\\NoAnalysisError\\DownloadError} \\
\hline
\rowcolor{lavender}
\multicolumn{4}{|p{\dimexpr\textwidth-2\tabcolsep}|}{\textbf{Description:} Saves overall analysis (generated by Generator) to User Activity table.} \\
\hline
\rowcolor{white}
save\_overall\_analysis & \makecell[lt]{\textit{user\_acc\_p}\\\textit{date}} & \textit{recent\_analysis} & \makecell[lt]{UserNotFoundError\\NoAnalysisError\\DownloadError} \\
\hline

\end{tabularx}
\caption{Exported Routines}

\end{table}
\end{comment}

%%%%%%%%%%%%%%%%%%%%%%%%%%%%%%%%%%%%%%%%%%%%%%%%%%%%%%%%%%%%%%%%%%%%%%


%%%%%%%%%%%%%%%%%%%%%%%%%%%%%%%%%%%%%%%%%%%%%%%%%%%%%%%%%%%%%%%%
\subsection{Module - Analysis Control Flow}
This module manages the video capture pipeline. It is responsible for validating appropriateness of the environment and 
quality assurance before analysis. Frames are processed directly for transmission. 

\subsubsection{Uses}
\noindent \textbf{\underline{\textit{server.py - Functions}}}\\
\begin{description}[style=nextline, leftmargin=2.5cm]
\item[\texttt{Flask}] 
    From "flask". Used to create the Flask web application instance and define API routes.

\item[\texttt{jsonify}] 
    From "flask". Used to convert Python dictionaries to JSON responses for API endpoints.

\item[\texttt{request}] 
    From "flask". Used to parse incoming HTTP request data, including JSON payloads and form data.

\item[\texttt{CORS}] 
    From "flask\_cors". Used to enable Cross-Origin Resource Sharing, allowing the mobile app to communicate with the backend.

\item[\texttt{threading.Lock}] 
    From Python standard library "threading". Used to create a lock for thread-safe access to the pose backend state.

\item[\texttt{base64.b64decode}] 
    From Python standard library "base64". Used to decode base64-encoded image strings sent from the mobile app.

\item[\texttt{np.frombuffer}] 
    From "numpy". Used to convert decoded image bytes into a NumPy array for image processing.

\item[\texttt{cv.imdecode}] 
    From "cv2" (OpenCV). Used to decode a NumPy array into an OpenCV image format (BGR).

\item[\texttt{cv.cvtColor}] 
    From "cv2" (OpenCV). Used to convert images between color spaces (BGR to GRAY) for brightness calculation.

\item[\texttt{cv.resize}] 
    From "cv2" (OpenCV). Used to resize frames to a standard dimension (640x480) before processing.

\item[\texttt{cv.flip}] 
    From "cv2" (OpenCV). Used to mirror images horizontally when using the front camera.

\item[\texttt{os.path.dirname}] 
    From Python standard library "os". Used to get the directory path of the current file.

\item[\texttt{os.path.abspath}] 
    From Python standard library "os". Used to get the absolute path of the current file.

\item[\texttt{os.path.join}] 
    From Python standard library "os". Used to construct the full file path to the pose landmarker model.

\item[\texttt{os.environ}] 
    From Python standard library "os". Used to set environment variables (OPENCV\_VIDEOIO\_MSMF\_ENABLE\_HW\_TRANSFORMS).

\item[\texttt{exists}] 
    From "os.path". Used to check if the pose landmarker model file exists at the specified path.

\item[\texttt{RuntimeError}] 
    Python built-in exception. Used to raise an error when the model file is missing.

\item[\texttt{threading}] 
    Python standard library module. Used for thread management and locks.
\end{description}

\noindent \textbf{\underline{\textit{server.py - Variables}}}\\
\begin{description}[style=nextline, leftmargin=2.5cm]
\item[\texttt{app}] 
    Flask application instance created for running the server.

\item[\texttt{pose\_app}] 
    Global instance of the PoseBackend class that manages pose detection state.

\item[\texttt{here}] 
    Variable storing the directory path of the current Python file.

\item[\texttt{model\_path}] 
    Variable storing the full file system path to the pose\_landmarker\_heavy.task model file.
\end{description}

\noindent \textbf{\underline{\textit{server.py - Types}}}\\
\begin{description}[style=nextline, leftmargin=2.5cm]
\item[\texttt{None}] 
    No external types are explicitly imported. Type hints use built-in types (str, bool, float, int) and OpenCV/NumPy types implicitly.
\end{description}
\noindent \textbf{\textit{\underline{poseService.ts}}}

\noindent \textbf{\underline{\textit{poseService.ts - Functions}}}\\
\begin{description}[style=nextline, leftmargin=2.5cm]
\item[\texttt{axios.post}] 
    From "axios". Used to send HTTP POST requests to the backend server for frame processing, pre-checking, and resetting the backend state.
\end{description}

\noindent \textbf{\underline{\textit{poseService.ts - Variables}}}\\
\begin{description}[style=nextline, leftmargin=2.5cm]
\item[\texttt{SERVER\_URL}] 
    Constant string variable storing the backend server address and port (http$://10.0.0.34:5001$). Used as the base URL for all API calls to the Flask server.

\item[\texttt{mirrored}] 
    Local boolean variable within processFrame and precheckFrame functions. Set to true when facing is "front", indicating the image needs to be mirrored.

\item[\texttt{res}] 
    Local variable storing the axios response object containing the backend response data.

\item[\texttt{headers}] 
    Configuration object passed to axios.post specifying Content-Type as "application\/json".

\item[\texttt{timeout}] 
    Configuration value set to 15000ms for processFrame and resetBackend, 10000ms for precheckFrame.
\end{description}

\noindent \textbf{\underline{\textit{poseService.ts - Types}}}\\
\begin{description}[style=nextline, leftmargin=2.5cm]
\item[\texttt{Side}] 
    Type union defining valid knee sides: "RIGHT" | "LEFT". Used to specify which knee is being analyzed.

\item[\texttt{Facing}] 
    Type union defining camera orientation: "front" | "back". Used to determine if image mirroring is needed.

\item[\texttt{Landmark}] 
    Type defining a 3D pose landmark with x, y, z coordinates (all numbers). Represents a detected body point.

\item[\texttt{Connection}] 
    Type defining a tuple [number, number] representing an edge between two landmark indices. Used to draw skeleton connections.

\item[\texttt{ProcessFrameResponse}] 
    Type defining the response structure from the \/process\_frame endpoint, containing landmarks, connections, metrics, and optional error.

\item[\texttt{PrecheckFrameResponse}] 
    Type defining the response structure from the \/precheck\_frame endpoint, containing ok status, lighting and visibility flags, and message.
\end{description}

\subsubsection{Syntax}

\textbf{Exported Constants}\\
\noindent \underline{\textit{poseService.ts}}\\
There are no exported constants.\\
\noindent \underline{\textit{server.py}}\\
No relevant exported constants\\


\noindent \textbf{Exported Access Programs}\\
\noindent \underline{\textit{server.py}}\\
\begin{longtable}{
|>{\columncolor{softgray}}p{4.2cm}|
>{\columncolor{white}}p{4.2cm}|
>{\columncolor{white}}p{4.2cm}|
}
\hline
\rowcolor{headerblue}
\textcolor{white}{\textbf{Input}} & 
\textcolor{white}{\textbf{Output}} & 
\textcolor{white}{\textbf{Exceptions}} \\
\hline
\endfirsthead

\hline
\rowcolor{headerblue}
\textcolor{white}{\textbf{Input}} & 
\textcolor{white}{\textbf{Output}} & 
\textcolor{white}{\textbf{Exceptions}} \\
\hline
\endhead

\multicolumn{3}{c}{\textit{Continued on next page}} \\
\endfoot
\endlastfoot

% reset_backend
\multicolumn{3}{|>{\columncolor{lightblue}}p{13.5cm}|}{\textbf{reset\_backend():} Flask route handler for POST /reset endpoint. Resets pose analyzer state and restarts calibration.} \\
\hline
\rowcolor{white}
\texttt{None} &
\texttt{JSON} - Returns \texttt{\{ok: true\}} &
\texttt{Exception} - Throws exception if reset fails \\
\hline

% precheck_frame
\multicolumn{3}{|>{\columncolor{lightblue}}p{13.5cm}|}{\textbf{precheck\_frame():} Flask route handler for POST /precheck\_frame endpoint. Validates lighting, frame composition, and knee visibility.} \\
\hline
\rowcolor{white}
\texttt{JSON body:} \texttt{imageBase64: string}, \texttt{side: string}, \texttt{mirrored: boolean}, \texttt{legsOnly: boolean} &
\texttt{JSON} - Returns \texttt{\{ok, tooDark, inFrame, kneeVisible, message\}} or error response &
\texttt{Exception} - Throws exception if frame decoding fails or processing error occurs \\
\hline

% process_frame
\multicolumn{3}{|>{\columncolor{lightblue}}p{13.5cm}|}{\textbf{process\_frame():} Flask route handler for POST /process\_frame endpoint. Processes frame for pose detection and metric calculation.} \\
\hline
\rowcolor{white}
\texttt{JSON body:} \texttt{imageBase64: string}, \texttt{side: string}, \texttt{mirrored: boolean}, \texttt{legsOnly: boolean} &
\texttt{JSON} - Returns \texttt{\{landmarks, connections, metrics\}} or error response with default metrics &
\texttt{Exception} - Throws exception if frame decoding fails or processing error occurs \\
\hline

\caption{Exported Access Programs - server.py} \label{sec:serverExports}
\end{longtable}

%%%%%%%%%%%%%%%%%%%%%%%%%%%%%%%%%%%%%%%%%%%%%%%
\noindent \underline{\textit{poseService.ts}}\\
\begin{longtable}{
|>{\columncolor{softgray}}p{4.2cm}|
>{\columncolor{white}}p{4.2cm}|
>{\columncolor{white}}p{4.2cm}|
}
\hline
\rowcolor{headerblue}
\textcolor{white}{\textbf{Input}} & 
\textcolor{white}{\textbf{Output}} & 
\textcolor{white}{\textbf{Exceptions}} \\
\hline
\endfirsthead

\hline
\rowcolor{headerblue}
\textcolor{white}{\textbf{Input}} & 
\textcolor{white}{\textbf{Output}} & 
\textcolor{white}{\textbf{Exceptions}} \\
\hline
\endhead

\multicolumn{3}{c}{\textit{Continued on next page}} \\
\endfoot
\endlastfoot

% processFrame
\multicolumn{3}{|>{\columncolor{lightblue}}p{13.5cm}|}{\textbf{processFrame(imageBase64, side, facing):} Sends a single video frame to the backend for pose detection and metric calculation.} \\
\hline
\rowcolor{white}
\texttt{imageBase64: string} - Base64 encoded image frame, \texttt{side: Side} - "RIGHT" or "LEFT", \texttt{facing: Facing} - "front" or "back" &
\texttt{Promise<ProcessFrameResponse>} - Returns landmarks, connections, and metrics &
\texttt{Error} - Throws error if request fails or timeout occurs (15000ms) \\
\hline

% precheckFrame
\multicolumn{3}{|>{\columncolor{lightblue}}p{13.5cm}|}{\textbf{precheckFrame(imageBase64, side, facing):} Performs pre-recording environment check including lighting and knee visibility.} \\
\hline
\rowcolor{white}
\texttt{imageBase64: string} - Base64 encoded image frame, \texttt{side: Side} - "RIGHT" or "LEFT", \texttt{facing: Facing} - "front" or "back" &
\texttt{Promise<PrecheckFrameResponse>} - Returns ok status, lighting flags, visibility flags, and message &
\texttt{Error} - Throws error if request fails or timeout occurs (10000ms) \\
\hline

% resetBackend
\multicolumn{3}{|>{\columncolor{lightblue}}p{13.5cm}|}{\textbf{resetBackend():} Resets the backend pose analyzer state and restarts calibration.} \\
\hline
\rowcolor{white}
\texttt{None} &
\texttt{Promise<void>} - Resolves when reset is complete &
\texttt{Error} - Throws error if request fails or timeout occurs (15000ms) \\
\hline

\caption{Exported Access Programs - poseService.ts} \label{sec:poseServiceExports}
\end{longtable}
\subsubsection{Semantics}

\noindent \textbf{State Variables} \\
\textit{F}: set of video frames\\\
\textit{C}: calibrating \\
\textit{cal\_time\_left}: time left for calibration\\
\textit{P}: frame responses set \\
\textit{L}: set of landmark arrays\\
\textit{url}: set of backend urls \\
\textit{$T_{process}$}: processFrame timeout value = 15000 ms\\
\textit{$T_{pre_check}$}: precheckFrame timeout value = 10000 ms\\
\textit{$T_{reset}$}: timeout value for resetBackend() = 15000 ms\\ 


\noindent \textbf{State Invariant} \\
\noindent
$\forall$ $b \in \text{Base64}: \text{b64\_to\_bgr\_image}(b) \in F$\\
$\forall$ $\text{operations on pose\_app}: \text{lock must be acquired first}$\\
$resetBackend() \implies analyzer = new Analyzer() \land c = true \cap cal\_time\_left = 10.0$
$\forall response \in P: \text{response.landmarks} \in L | response.connections \in C$\\
$true \implies T_{process} = 15000 \cap T_{precheck} = 10000 \cap T_{reset} = 15000$ \\
$processFrame endpoint = SERVER\_URL + /process\_frame$\\
$precheckFrame endpoint = SERVER\_URL + /precheck\_frame$\\
$resetBackend endpoint = SERVER\_URL + /reset$\\


\subsubsection{Environment variables}
NA\\

\subsubsection{Assumptions}
\begin{itemize}
    \item All frames resize to 640x480 pixels.
    \item When \texttt{mirrored = true} (front camera), the frame is horizontally flipped. 
    This assumes the patient is facing the camera and wants natural left-right orientation.
    \item A single user's frames can be processed at a time (this falls under the overarching assumption that the
    application can be used by one person at a time only).
    \item Images are assumed to be base64-encoded JPEG or PNG format. No other formats are supported.
    \item Timeout values assume reasonable response times.
    \item Reset function is assumed to be called before re-recording or new recording.
\end{itemize}

\subsubsection{Access Routine Semantics - Exported}
\texttt{f}: set of video frames\\\
\texttt{c}: calibrating \\
\texttt{cal\_time\_left}: time left for calibration\\
\textit{p}: frame responses set \\
\textit{l}: set of landmark arrays\\
\textit{url}: set of backend urls \\
\textit{b}: set of base64 encoded frames\\
\texttt{s}: {'RIGHT' | 'LEFT'}\\
\texttt{side}: $\in s$
\texttt{facing} : {'front' | 'back'}\\
\texttt{response}: PrecheckFrameResponse \\
\textit{J}: set of all JSON strings\\
\texttt{pose\_app}: PoseBackend()\\ 
\newline
\textbf{poseService.ts}\\
\newline
\textit{processFrame}:\\
\textbf{Pre-condition}:$\forall j \in J: j.imageBase64 \in b \cap 
j.side \in s \cap j.mirrored \in \{true, false\} \cap cameraReady = true \cap pose\_app.analyzer \neq \emptyset$ \\
\textbf{Function}: processFrame(fr, side, facing) \\
\textbf{Post-condition}:
$fr \in f \implies pose\_app.analyzer \neq \emptyset
\cap \{rep\_count\}(t_2) \geq \{rep\_count\}(t_1)
\forall t_2 > t_1, rom\_degree = max\_degree - min\_degree \geq 0
\cap angle \in [0, 180] \cap c(t_2) = false \implies c(t_1) = false 
\cup cal\_time\_left(t_1) \leq 0$
\textit{precheckFrame}:\\
\textbf{Pre-condition}: $\forall j \in J: j.imageBase64 \in b \cap
j.side \in s \cap j.mirrored \in {true, false} \cap cameraReady = true \cap 
pose\_app.analyzer \neq \emptyset$ \\
\textbf{Function}: precheckFrame(fr, s, facing)\\
\textbf{Post-condition}: $\forall j \in J: response.ok \in \{true, false\} \cap response.message \neq \emptyset:
(response.ok = true \implies response.tooDark = false \cap response.inFrame = true \cap response.kneeVisible = true)$
$\cup$ $\text{response.ok = false} \implies \text{response.message describes failure reason}$ \\
\newline
\textit{resetBackend()}: \\
\textbf{Pre-condition}: $pose\_app \neq \emptyset$ \\
\textbf{Function}: resetBackend()\\
\textbf{Post-condition}: $pose\_app.analyzer = new Analyzer() \cap c = true \cap cal\_time\_left = 10 s
\cap rep\_count = 0 \cap min\_degree, max\_degree, rom\_degree = 0$\\
\newline
\noindent \textbf{server.py}\\
\newline
\textit{reset\_backend (POST/reset endpoint)} \\
\textbf{Pre-condition:} $pose\_app \neq \emptyset$ \\
\textbf{Function:} reset\_backend() \\
\textbf{Post-condition:} $ pose\_app.analyzer = new Analyzer()$ \\
$ is\_calibrating = true$\\
$ cal\_time\_left = 10.0$\\
$rep\_count = 0$\\

\noindent \textit{precheck\_frame (POST/precheck\_frame endpoint)}\\
\textbf{Pre-condition}: $\forall j \in J: j.imageBase64 \in b \cap
j.side \in s \cap j.mirrored \in {true, false} \cap cameraReady = true$ \\
\textbf{Function}: precheckFrame() \\
\textbf{Post-condition}: $response.message \neq \emptyset:
(response.ok = true \implies response.tooDark = false \cap response.inFrame = true \cap response.kneeVisible = true)$
$\cup$ $\text{response.ok = false} \implies \text{response.message describes failure reason}$ \\

\noindent \textit{process\_frame(POST/process\_frame endpoint)}\\
\textbf{Pre-condition}:$\forall j \in J: j.imageBase64 \in b \cap 
j.side \in s \cap j.mirrored \in \{true, false\} \cap cameraReady = true \cap pose\_app.analyzer \neq \emptyset$ \\
\textbf{Function}: processFrame()\\
\textbf{Post-condition}: $\{rep\_count\}(t_2) \geq \{rep\_count\}(t_1)
\forall t_2 > t_1, rom\_degree = max\_degree - min\_degree \geq 0
\cap angle \in [0, 180] \cap c(t_2) = false \implies c(t_1) = false 
\cup cal\_time\_left(t_1) \leq 0$ \\



\subsubsection{Local Functions}
\subsubsection{Access Routine Semantics - Local}

%%%%%%%%%%%%%%%%%%%%%%%%%%%%%%%%%%%%%%%%%%%%%%%%%%%%%%%%%%%%%%%%%%%%
\subsection{Module - Analyze}

\subsubsection{Uses}
%Relevant files: analyze.py, server.py (backend)
\noindent \textbf{\underline{\textit{analyze.py - Variables}}}\\
\begin{description}[style=nextline, leftmargin=2.5cm]
\item[\texttt{min\_angle}] 
    Stores the minimum knee angle observed during calibration of range of motion.

\item[\texttt{max\_angle}] 
    Stores the maximum knee angle observed during calibration of range of motion.

\item[\texttt{rep\_count}] 
    Tracks the number of completed repetitions.

\item[\texttt{rep\_state}] 
    Current phase of repetition ("Extension", "Flexion").

\item[\texttt{cue\_state}] 
    Current feedback cue state ("CALIBRATING", "GOOD\_EXTENSION", \\"GOOD\_FLEXION", "GETTING\_READY", "OUT\_OF\_FRAME").

\item[\texttt{smoothed\_angle}] 
    Exponentially smoothed angle value to reduce jitter.

\item[\texttt{cycle\_min\_angle}] 
    Minimum angle recorded within the current repetition cycle.

\item[\texttt{cycle\_max\_angle}] 
    Maximum angle recorded within the current repetition cycle.

\item[\texttt{last\_rep\_time}] 
    Timestamp of the last completed repetition.
\end{description}

\noindent \textbf{\underline{\textit{analyze.py - Types}}}\\
None\\

\noindent \textbf{\underline{\textit{analyze.py - Functions}}}\\
\begin{description}[style=nextline, leftmargin=2.5cm]
\item[\texttt{\_smooth\_angle(angle)}] 
    Applies smoothing to angle measurements/calculates acceptable angle values.

\item[\texttt{\_calibrate\_with\_angle(angle)}] 
    Updates min/max angle values during calibration window.

\item[\texttt{\_rep\_update(angle)}] 
    Core repetition detection logic using zone-based thresholds.

\item[\texttt{start\_calibration(duration\_s)}] 
    Begins calibration phase for specified duration in seconds.

\item[\texttt{update(angle)}] 
    Main entry point for processing each new angle measurement.

\item[\texttt{calc\_rom()}] 
    Returns range of motion as max\_angle minus min\_angle.

\item[\texttt{summary()}] 
    Returns dictionary containing all current analysis metrics.
\end{description}

\noindent \textbf{\underline{\textit{server.py - Variables}}}\\
\begin{description}[style=nextline, leftmargin=2.5cm]
\item[\texttt{pose\_app}] 
    Global singleton managing analyzer state and camera.

\item[\texttt{analyzer}] 
    Instance of Analyzer class accessed via pose\_app.analyzer.

\item[\texttt{lock}] 
    Ensures thread-safe access to analyzer during concurrent requests.
\end{description}

\noindent \textbf{\underline{\textit{server.py - Types}}}\\
None\\

\noindent \textbf{\underline{\textit{server.py - Functions}}}\\
\begin{description}[style=nextline, leftmargin=2.5cm]
\item[\texttt{Flask}] 
    Creates the Flask web application instance and defines API routes.

\item[\texttt{jsonify}] 
    Converts Python dictionaries to JSON HTTP responses.

\item[\texttt{request}] 
    Accesses incoming request data including JSON body.

\item[\texttt{b64\_to\_bgr\_image(b64\_str)}] 
    Decodes base64 string to OpenCV BGR image format.

\item[\texttt{process\_frame()}] 
    Receives frame, calculates knee angle, updates analyzer, returns metrics.

\item[\texttt{check\_side\_visibility(landmarks, side)}] 
    Determines if specified knee and adjacent joints are visible in frame.

\item[\texttt{Thread}] 
    Enables concurrent request handling via Flask's threaded=True.
\end{description}

\subsubsection{Syntax}

\textbf{Exported Constants} \\
\noindent \textbf{\underline{\textit{analyze.py, server.py}}}\\
None\\

\noindent \textbf{Exported Access Programs}\\
\noindent \underline{\textit{\textbf{analyze.py}}}\\
\begin{longtable}{
|>{\columncolor{softgray}}p{4.2cm}|
>{\columncolor{white}}p{4.2cm}|
>{\columncolor{white}}p{4.2cm}|
}
\hline
\rowcolor{headerblue}
\textcolor{white}{\textbf{Input}} & 
\textcolor{white}{\textbf{Output}} & 
\textcolor{white}{\textbf{Exceptions}} \\
\hline
\endfirsthead

\hline
\rowcolor{headerblue}
\textcolor{white}{\textbf{Input}} & 
\textcolor{white}{\textbf{Output}} & 
\textcolor{white}{\textbf{Exceptions}} \\
\hline
\endhead

\multicolumn{3}{c}{\textit{Continued on next page}} \\
\endfoot
\endlastfoot

% start_calibration
\multicolumn{3}{|>{\columncolor{lightblue}}p{13.5cm}|}{\textbf{start\_calibration(duration\_s):} Begins calibration phase to record min and max knee angles.} \\
\hline
\rowcolor{white}
\texttt{duration\_s: float} &
\texttt{None} &
\texttt{None} \\
\hline

% update
\multicolumn{3}{|>{\columncolor{lightblue}}p{13.5cm}|}{\textbf{update(angle):} Main entry point for processing each new angle measurement. Updates smoothed angle, calibration state, and rep detection.} \\
\hline
\rowcolor{white}
\texttt{angle: float} &
\texttt{None} &
\texttt{angle: float} \\
\hline

% calc_rom
\multicolumn{3}{|>{\columncolor{lightblue}}p{13.5cm}|}{\textbf{calc\_rom():} Calculates and returns the current range of motion.} \\
\hline
\rowcolor{white}
\texttt{None} &
\texttt{None} &
\texttt{float} - Range of motion in degrees \\
\hline

% summary
\multicolumn{3}{|>{\columncolor{lightblue}}p{13.5cm}|}{\textbf{summary():} Returns a dictionary containing all current analysis metrics.} \\
\hline
\rowcolor{white}
\texttt{None} &
\texttt{dict} - Contains min\_degree, max\_degree, rom\_degree, calibrating, cal\_time\_left, rep\_count, current\_rep\_duration, avg\_rep\_duration, rep\_state, cue\_state &
\texttt{None} \\
\hline
\caption{Access Routines for Analysis Module}
\label{table:access_routines_analysis}
\end{longtable}
\noindent \underline{\textit{\textbf{server.py}}}\\
None 

\subsubsection{Semantics}

\noindent \textbf{State Variables} \\
\textit{S}: set containing "Extension" and "Flexion"\\
\texttt{ex\_state}: Exercise state\\
\textit{C}: {"CALIBRATING", "GOOD\_EXTENSION", "GOOD\_FLEXION", "GETTING\_READY", "OUT\_OF\_FRAME"} \\
\texttt{started}: (True, False) indicating whether exercise has started.\\
\texttt{angle}: angle measurement of landmark vectors \\
\texttt{min\_angle}: min\_angle\\
\texttt{max\_angle}: max\_angle\\
\texttt{is\_calibrating}: (True, False) indicating whether calibration is taking place.\\
\texttt{cue}: exercise feedback string \\
\textit{A}: set of Analyzer instances \\



\noindent \textbf{State Invariants} \\
$\forall$ $angle \in \{angles\}, angle \in 0 \leq angle \leq 180$ \\
$ (\texttt{started} = True \implies \texttt{is\_calibrating} = False \cup (\texttt{ex\_state} \in S$ $|$ $ex\_state = 
"Flexion" \cup "Extension")) 
\cup (\texttt{started} = False \implies (is\_calibrating = True))$\\
$\forall \texttt{ex\_state} \in S, \texttt{cue} \in C, ex\_state = "Extension" \implies \texttt{cue} = "GOOD EXTENSION"
\cup ex\_state = "Flexion" \implies \texttt{cue} = "GOOD FLEXION"$\\
$\forall a \in A, |a| = 1$\\

\subsubsection{Environment Variables}
NA\\

\subsubsection{Assumptions}
\begin{itemize}
    \item Basic process threading rules apply when using an Analyzer instance. Only one process can function properly.
\end{itemize}

\subsubsection{Access Routine Semantics - Exported}

\subsubsection{Local Functions}

\subsubsection{Access Routine Sematics - Local}

%%%%%%%%%%%%%%%%%%%%%%%%%%%%%%%%%%%%%%%%%%%%%%%%%%%%%%
\subsection{Module - Feedback}
This module's role is to generate real-time feedback based on form analysis of the exercise that is currently being performed.
It provides guidance based on form calculations and alerts the user on screen. 

%FeedbackOverlay.ts

\subsubsection{Uses}
\noindent \underline{\textbf{\textit{FeedbackOverlay.ts - Variables}}}
\begin{description}

\item[\texttt{normalized}] 
    Computed normalized angle value derived from metrics, used to determine feedback threshold comparisons.

\item[\texttt{FLEXION\_THRESHOLD}] 
    Constant value 0.7. Threshold for detecting when knee is sufficiently bent (normalized angle below 0.3 triggers flexion feedback).

\item[\texttt{EXTENSION\_THRESHOLD}] 
    Constant value 0.7. Threshold for detecting when knee is sufficiently extended (normalized angle above 0.7 triggers extension feedback).

\item[\texttt{CAL\_DURATION\_S}] 
    Constant value 10.0. Duration of calibration phase in seconds.
\end{description}
\noindent \underline{\textbf{\textit{analyze.py - Variables}}}
\begin{description}
\item[\texttt{cue\_state}] 
    String variable indicating current feedback cue ("CALIBRATING", "GOOD\_EXTENSION", "GOOD\_FLEXION", "GETTING\_READY", "OUT\_OF\_FRAME"). Generates user-readable messages.

\item[\texttt{rep\_state}] 
    String variable indicating current repetition phase ("Extension", "Flexion", "Ready"). Used to determine appropriate guidance message.

\item[\texttt{rep\_count}] 
    Integer tracking number of completed repetitions. Displayed in feedback subtitle.

\item[\texttt{calibrating}] 
    Boolean indicating whether calibration phase is active. Determines feedback type.

\item[\texttt{cal\_time\_left}] 
    Float indicating seconds remaining in calibration. Displayed in calibration feedback.

\item[\texttt{min\_angle}, \texttt{max\_angle}] 
    Float values representing calibrated range of motion. Used to calculate normalized angle for threshold comparisons.
\end{description}

\noindent \underline{\textbf{\textit{FeedbackOverlay.ts - Types}}}
\begin{description}
\item[\texttt{Metrics}] 
    Type definition representing exercise metrics structure containing angle, min\_degree, max\_degree, rom\_degree, rep\_count, rep\_state, calibrating, cal\_time\_left, current\_rep\_duration, avg\_rep\_duration.

\item[\texttt{FeedbackInfo}] 
    Type definition representing feedback output structure with fields: type ("error" | "warning" | "info" | "success"), icon (string), title (string), subtitle (string), color (string).
\end{description}
\noindent \underline{\textbf{\textit{analyze.py - Functions}}}
\begin{description}
\item[\texttt{summary()}] 
    Returns dictionary containing all current metrics including cue\_state, rep\_state, rep\_count, calibrating, cal\_time\_left, min\_angle, max\_angle, rom\_degree.

\item[\texttt{update(angle)}] 
    Processes angle measurements and updates internal state including cue\_state.
\end{description}
\noindent \underline{\textbf{\textit{FeedbackOverlay.ts - Functions}}}
\begin{description}
\item[\texttt{getFeedback(metrics)}] 
    Core feedback generation function. Evaluates calibration state, rep state, and normalized angle to return appropriate feedback message, icon, and color.

\item[\texttt{getNormalizedAngle(metrics)}] 
    Helper function that converts raw angle measurement to a normalized value between 0 and 1 based on the user's calibrated min and max range of motion.
\end{description}

\subsubsection{Syntax}
\textbf{Exported Constants} \\
None\\

\noindent \textbf{Exported Access Programs}\\

\noindent \underline{\textbf{\textit{FeedbackOverlay.ts}}}\\
None\\
\noindent \underline{\textbf{\textit{analyze.py}}}\\
\begin{longtable}{
|>{\columncolor{softgray}}p{4.2cm}|
>{\columncolor{white}}p{4.2cm}|
>{\columncolor{white}}p{4.2cm}|
}
\hline
\rowcolor{headerblue}
\textcolor{white}{\textbf{Input}} & 
\textcolor{white}{\textbf{Output}} & 
\textcolor{white}{\textbf{Exceptions}} \\
\hline
\endfirsthead

\multicolumn{3}{|>{\columncolor{lightblue}}p{13.5cm}|}{\textbf{summary():} Returns dictionary containing cue\_state, rep\_state, rep\_count, calibrating, cal\_time\_left, min\_angle, max\_angle, rom\_degree, current\_rep\_duration, avg\_rep\_duration.} \\
\hline
\rowcolor{white}
\texttt{None} &
\texttt{dict} &
\texttt{None} \\
\hline
\end{longtable}

\subsubsection{Sematics}
\textbf{State Variables}\\
\subsubsection{Environment Variables}
\subsubsection{Assumptions}
\subsubsection{Access Routine Semantics - Exported}
\subsubsection{Local Functions}
\subsubsection{Access Routine Semantics - Local}

%%%%%%%%%%%%%%%%%%%%%%%%%%%%%%%%%%%%%%%%%%%%%%%%%%%%%
\subsection{Module - Metrics}
\begin{comment}
This module measures metrics (\textit{angle, duration, height}) that are necessary to construct any corrections to the user's form. \\

TODO

\subsubsection{Uses}

\subsubsection{Syntax}

\textbf{Exported Constants}

\noindent \textbf{Exported Access Programs}
\begin{table}[h]
\centering
\setlength{\extrarowheight}{3pt}
\setlength{\arrayrulewidth}{1.2pt}
\begin{tabular}{|l|l|l|l|}
\hline
\rowcolor{darkpurple}
\color{white}\textbf{Access Routine Name} & \color{white}\textbf{Input} & \color{white}\textbf{Output} & \color{white}\textbf{Exceptions} \\
\hline


\rowcolor{lavender}
\multicolumn{4}{|p{\dimexpr\textwidth-2\tabcolsep}|}{\textbf{Description:} obtains a frame from Main module} \\
\hline
\rowcolor{white}
get\_frame & \makecell[lt]{\textit{processed\_frame}} & \makecell[lt]{NA} & \makecell[lt]{NoFrameError} \\
\hline

\rowcolor{lavender}
\multicolumn{4}{|p{\dimexpr\textwidth-2\tabcolsep}|}{\textbf{Description:} calculates metrics from frame provided by Main module} \\
\hline
\rowcolor{white}
calculate\_frame & \makecell[lt]{\textit{processed\_frame}} & \makecell[lt]{\textit{angle}\\\textit{height}\\\textit{duration}} & \makecell[lt]{NoFrameError} \\
\hline

\rowcolor{lavender}
\multicolumn{4}{|p{\dimexpr\textwidth-2\tabcolsep}|}{\textbf{Description:} obtains metrics that need to be adjusted by the user (calls Analysis)} \\
\hline
\rowcolor{white}
fixed\_frame & \makecell[lt]{\textit{NA}} & \makecell[lt]{\textit{fixed\_metrics}} & \makecell[lt]{NoFrameError} \\
\hline



\end{tabular}
\caption{Access Routines for Analysis Module}
\label{table:access_routines}
\end{table}

\subsubsection{Semantics}

\noindent \textbf{State Variables} \\


%Check requirements doc
\noindent \textbf{State Invariant} \\


%Environment variables?

\subsubsection{Assumptions}

\subsubsection{Access Routine Semantics}

%\noindent \wss{accessProg}():
%\begin{itemize}
%\item transition: \wss{if appropriate} 
%\item output: \wss{if appropriate} 
%\item exception: \wss{if appropriate} 
%\end{itemize}

\subsubsection{Local Functions}
\begin{table}[h]
\centering
\setlength{\extrarowheight}{3pt}
\setlength{\arrayrulewidth}{1.2pt}
\begin{tabular}{|l|l|l|l|}
\hline
\rowcolor{darkpurple}
\color{white}\textbf{Access Routine Name} & \color{white}\textbf{Input} & \color{white}\textbf{Output} & \color{white}\textbf{Exceptions} \\
\hline

% measures the angle between the starting and final plane
\rowcolor{lavender}
\multicolumn{4}{|p{\dimexpr\textwidth-2\tabcolsep}|}{\textbf{Description:} measures the angle between the starting and final plane} \\
\hline
\rowcolor{white}
measure\_angle & \makecell[lt]{\textit{initial\_plane}\\\textit{final\_plane}} & \makecell[lt]{\textit{angle}} & \makecell[lt]{FrameEmptyError} \\
\hline

% measures difference in initial and final height
\rowcolor{lavender}
\multicolumn{4}{|p{\dimexpr\textwidth-2\tabcolsep}|}{\textbf{Description:} measures difference in initial and final height} \\
\hline
\rowcolor{white}
measure\_height & \makecell[lt]{\textit{initial\_height}\\\textit{final\_height}} & \makecell[lt]{\textit{height}} & \makecell[lt]{FrameEmptyError} \\
\hline

% measures duration of the performed exercise
\rowcolor{lavender}
\multicolumn{4}{|p{\dimexpr\textwidth-2\tabcolsep}|}{\textbf{Description:} measures duration of the performed exercise} \\
\hline
\rowcolor{white}
measure\_duration & \makecell[lt]{\textit{initial\_time}\\\textit{final\_time}} & \makecell[lt]{\textit{duration}} & \makecell[lt]{FrameEmptyError} \\
\hline

\end{tabular}
\caption{Metrics Access Routines}
\label{tab:metrics_routines}
\end{table}

\subsubsection{Local Functions}
\end{comment}
 %%%%%%%%%%%%%%%%%%%%%%%%%%%%%%%%%%%%%%%%%%%%%%%%%%%%%%%%%%%%%%%%
 \subsection{Module - Draw}
\begin{comment}
 This module is responsible for drawing landmarks by using \textit{MediaPipe} (external library).


\subsubsection{Uses}
\begin{itemize}
  \item starting\_frame
  \item processed\_frames\_list
  \item unprocessed\_frames\_list
  \item feedback\_instructions

\end{itemize}

\newpage
\subsubsection{Syntax}

\textbf{Exported Constants}
\begin{itemize}
  \item \textit{can\_record}: good\_to\_record
  \item \textit{fixing\_instructions}: feedback\_instructions
  \item \textit{processed}: processed\_frames\_list
\end{itemize}


\noindent \textbf{Exported Access Programs}
\begin{table}[h]
\centering
\setlength{\extrarowheight}{3pt}
\setlength{\arrayrulewidth}{1.2pt}
\begin{tabular}{|l|l|l|l|}
\hline
\rowcolor{darkpurple}
\color{white}\textbf{Access Routine Name} & \color{white}\textbf{Input} & \color{white}\textbf{Output} & \color{white}\textbf{Exceptions} \\
\hline

% Checks if the recording environment is appropriate (e.g., lighting, background)
\rowcolor{lavender}
\multicolumn{4}{|p{\dimexpr\textwidth-2\tabcolsep}|}{\textbf{Description:} Checks if the recording environment is appropriate (e.g., lighting, background)} \\
\hline
\rowcolor{white}
record\_check & \makecell[lt]{\textit{starting\_frame}} & \makecell[lt]{\textit{good\_to\_record}} & \makecell[lt]{UserNotFound\\Error\\DownloadError} \\
\hline

% Analyzes the user's form against a correct form model (determined by Analysis)
\rowcolor{lavender}
\multicolumn{4}{|p{\dimexpr\textwidth-2\tabcolsep}|}{\textbf{Description:} Analyzes the user's starting form against a correct form model} \\
\hline
\rowcolor{white}
form\_check & \makecell[lt]{\textit{processed\_frames\_list}} & \makecell[lt]{\textit{feedback\_instructions}} & \makecell[lt]{UserNotFound\\Error\\DownloadError} \\
\hline

% Overlays pose estimation points onto the video feed for user feedback
\rowcolor{lavender}
\multicolumn{4}{|p{\dimexpr\textwidth-2\tabcolsep}|}{\textbf{Description:} Overlays pose estimation points onto the video feed for user feedback} \\
\hline
\rowcolor{white}
draw\_points & \makecell[lt]{\textit{unprocessed\_frames}} & \makecell[lt]{processed\_frames\_list} & \makecell[lt]{UserNotFound\\Error\\DownloadError} \\
\hline

\end{tabular}
\caption{Draw Access Routines}
\label{tab:video_routines}
\end{table}

\subsubsection{Semantics}

\noindent \textbf{State Variables} \\
\textit{frames}: unprocessed\_frames\_list\\


%Check requirements doc
\noindent \textbf{State Invariant} \\


%Environment variables?

\subsubsection{Assumptions}
\begin{itemize}
  \item User's complete form is in frame.
  \item The application will not draw points if the user does not comply with instructions given by \textit{record\_check}.
\end{itemize}
\subsubsection{Access Routine Semantics}
TODO


\subsubsection{Local Functions}
\begin{table}[h]
\centering
\setlength{\extrarowheight}{3pt}
\setlength{\arrayrulewidth}{1.2pt}
\begin{tabular}{|l|l|l|l|}
\hline
\rowcolor{darkpurple}
\color{white}\textbf{Access Routine Name} & \color{white}\textbf{Input} & \color{white}\textbf{Output} & \color{white}\textbf{Exceptions} \\
\hline
%MP call 
\rowcolor{lavender}
\multicolumn{4}{|p{\dimexpr\textwidth-2\tabcolsep}|}{\textbf{Description:} Calls MediaPipe's draw functionality} \\
\hline
\rowcolor{white}
mpipe\_draw & \makecell[lt]{\textit{unprocessed\_frames}} & \makecell[lt]{\textit{processed\_frames\_list}} & \makecell[lt]{MPbusyError} \\
\hline

\end{tabular}
\caption{Account Management Access Routines}
\label{tab:account_routines}
\end{table}
\end{comment}
%%%%%%%%%%%%%%%%%%%%%%%%%%%%%%%%%%%%%%%%%%%%%%%%%%%%%%%%%%%%%%%%%
\subsection{Module - UI} \label{sec:ui}
\begin{comment}
NA

%exercises.tsx


\subsubsection{Uses}

\subsubsection{Syntax}

\textbf{Exported Constants}

\noindent \textbf{Exported Access Programs}
\begin{longtable}{|p{0.25\textwidth}|p{0.25\textwidth}|p{0.25\textwidth}|p{0.25\textwidth}|}
\hline
\rowcolor{headerblue}
\color{white}\textbf{Access Routine Name} & \color{white}\textbf{Input} & \color{white}\textbf{Output} & \color{white}\textbf{Exceptions} \\
\hline
\endfirsthead

\hline
\rowcolor{headerblue}
\color{white}\textbf{Access Routine Name} & \color{white}\textbf{Input} & \color{white}\textbf{Output} & \color{white}\textbf{Exceptions} \\
\hline
\endhead

\hline
\endfoot

\rowcolors{2}{tableblue}{white}
\setlength{\extrarowheight}{3pt}
\setlength{\arrayrulewidth}{1.2pt}

placeholder & \makecell[lt]{\textit{NA}} & \makecell[lt]{UserAccount\_P\\OR\\UserAccount\_PT} & \makecell[lt]{User\_not\_found\\Download\_error} \\
\hline


\end{longtable}

\subsubsection{Semantics}
TBD \\
\noindent \textbf{State Variables} \\

TODO \\
%Check requirements doc
\noindent \textbf{State Invariant} \\
TODO \\

%Environment variables?

\subsubsection{Assumptions}

\subsubsection{Access Routine Semantics}

%\noindent \wss{accessProg}():
%\begin{itemize}
%\item transition: \wss{if appropriate} 
%\item output: \wss{if appropriate} 
%\item exception: \wss{if appropriate} 
%\end{itemize}

\subsubsection{Local Functions}
\begin{table}[h]
\centering
\rowcolors{2}{tableblue}{white}
\setlength{\extrarowheight}{3pt}
\setlength{\arrayrulewidth}{1.2pt}

\begin{tabularx}{\textwidth}{|>{\raggedright\arraybackslash}X|>{\raggedright\arraybackslash}X|>{\raggedright\arraybackslash}X|>{\raggedright\arraybackslash}X|}
\hline
\rowcolor{headerblue}
\color{white}\textbf{Access Routine Name} & \color{white}\textbf{Input} & \color{white}\textbf{Output} & \color{white}\textbf{Exceptions} \\

\hline
placeholder & \makecell[lt]{\textit{Name}\\\textit{Role}\\\textit{Birthday}\\\textit{License\_no/}\\\textit{Invite\_code}\\\textit{Email}\\\textit{Password}} & \makecell[lt]{UserAccount\_P\\or\\UserAccount\_PT} & \makecell[lt]{Input\_format\_invalid\\Required\_field\_empty\\Account\_already\_\\exists\\Incorrect\_creds\\Invite\_code\_expired} \\ %creates a profile
\hline


\end{tabularx}
\caption{UI Access Routines}
\label{tab:account_routines}
\end{table}



\subsubsection{Local Functions}

\end{comment}
%%%%%%%%%%%%%%%%%%%%%%%%%%%%%%%%%%%%%%%%%%%%%%%%%%%%%%%%%%%%%%%%%

\bibliographystyle {plainnat}
\bibliography {../../../refs/References}
\begin{comment}

\section{Appendix A}
For clarity and readability, the variables used throughout the document are provided below.

\begin{longtable}{|>{\raggedright\arraybackslash}p{0.22\textwidth}|>{\raggedright\arraybackslash}p{0.20\textwidth}|>{\raggedright\arraybackslash}p{0.38\textwidth}|>{\raggedright\arraybackslash}p{0.20\textwidth}|}
\caption{Variable used throughout MIS}\\
\hline
\rowcolor{headerblue}
\color{white}\textbf{Name} & \color{white}\textbf{Type} & \color{white}\textbf{Description} & \color{white}\textbf{Used in sections} \\
\hline
\endfirsthead

\caption[]{Variable used throughout MIS (continued)} \\
\hline
\rowcolor{headerblue}
\color{white}\textbf{Name} & \color{white}\textbf{Type} & \color{white}\textbf{Description} & \color{white}\textbf{Used in sections} \\
\hline
\endhead

\hline
\endfoot

\rowcolors{2}{tableblue}{white}
\setlength{\extrarowheight}{2pt}
\setlength{\arrayrulewidth}{1pt}

name & String & Name of the user & 7.1 \\
\hline
role & String & Values can be PT or patient & 7.1 \\
\hline
email & String & Used as login ID & 7.1 \\
\hline
password & String & Set by user for login purposes & 7.1 \\
\hline
birthday & String & User's birthday for verification & 7.1 \\
\hline
license\_no & Int & Verification of PT & 7.1 \\
\hline
invite\_code & String & Given by PT to patient for registration & 7.1 \\
\hline
is\_valid & Boolean & Used to check validity of user (invite\_code, license\_no) & ? \\
\hline
acc\_id & Int & Identifier assigned to app users & ? \\
\hline
acc\_double & Boolean & Indicates whether an account already exists in the system & ? \\
\hline
date & String & Date (YYYYMMDD) & 7.1 \\
\hline
time & String & Time (s) & 7.2 \\
\hline
recent\_analysis & ADT & Most recently generated analysis & 7.1 \\
\hline
recent\_report & Record & Generated report of recent\_video & TBD \\
\hline
recent\_video & String (JSON) & Last video recorded by patient & TBD \\
\hline
video & String (JSON) & A video in the user database & TBD \\
\hline
report & Record & Report(s) sent to Generator module for analysis OR requested by patient & TBD \\
\hline
ex\_instruct & String & Performance instructions of requested exercise & TBD \\
\hline
sets & String & Number of sets given by PT & 7.1 \\
\hline
repetitions & String & Repetitions of exercise per set & 7.1 \\
\hline
rest\_time & String & Rest time between sets & 7.1 \\
\hline
leg\_landmarks & Array & MediaPipe output of leg & 7.4 \\
\hline
processed\_frame & Image & Video frame with overlaid markers & 7.4 \\
\hline
PT\_comment & String & PT's comment on patient activity & 7.3 \\
\hline
overall\_analysis & Abstract object & the overall analysis of all patient activity & 7.4 \\ %check
\hline
processed\_frames\_list & List(Image object) & list of processed\_frames & 7.3 \\
\hline
feedback\_instructions & Abstract Object & instructions given from analysis to feedback module & 7.3 \\
\hline
fixed\_metrics & float & corrected metrics & 7.3 \\
\hline
initial/final\_\newline plane/time/height & float & initial and final measurements & 7.3 \\
\hline
%TODO
ins\_ui & String & PT's comment on patient activity & TBD \\
\hline
unprocessed\_frames & List(Image objects) & Unprocessed frames list & TBD \\
\hline
rt\_metrics & List(of float) & Real-time metrics & TBD \\
\hline
input\_frame & Image & Input frame & TBD \\
\hline
exercise & String & Selected exercise & TBD \\
\hline
insights & String & Generated insights & TBD \\
\hline
all\_reports & List[report] & All reports are stored in a list & TBD \\
\hline
starting\_frame & Image & The last frame before recording & TBD \\
\hline
good\_to\_record & String & A message informing user that the system is ready to record & TBD \\
\hline
\end{longtable}
\end{comment}

\section*{Appendix A}
\begin{longtable}{L R}
\caption{Database Schema}\\
\toprule
\rowcolor{headerblue}
\textcolor{white}{\textbf{Tables}} & \textcolor{white}{\textbf{Fields}} \\
\midrule
\endfirsthead

\multicolumn{2}{c}%
{{\bfseries \tablename\ \thetable{} -- continued from previous page}} \\
\toprule
\rowcolor{headerblue}
\textcolor{white}{\textbf{Tables}} & \textcolor{white}{\textbf{Fields}} \\
\midrule
\endhead

\bottomrule
\multicolumn{2}{r}{\textit{Continued on next page}} \\
\endfoot

\bottomrule
\endlastfoot

\textbf{activities} & \texttt{actid, createdAt, duration, email, exercise, timestamp, uid, analysis, avg\_rep\_duration, completed\_reps, completed\_sets, current\_angle, current\_rep\_duration, date\_performed, effort, max\_height, min\_height, name, pain, patient\_feedback, rep\_state, rom\_degree, satisfaction, target\_area, Database location} \\
\midrule
\textbf{comments} & \texttt{actid, author, cid, comment, date, uid} \\
\midrule
\textbf{exercises} & \texttt{description, reps, sets, title} \\
\midrule
\textbf{inviteCodes} & \texttt{active, physioId, used, usedBy} \\
\midrule
\textbf{userExercises} & \texttt{description, id, reps, sets, title, userid} \\
\midrule
\textbf{users} & \texttt{birthday, createdAt, email, inviteCode, name, physioId, role} \\
\midrule
\textbf{validLicenses} & \texttt{active} \\

\end{longtable}
\section*{Appendix --- Reflection}

The information in this section will be used to evaluate the team members on the
graduate attribute of Problem Analysis and Design.

\input{../../Reflection.tex}



\begin{enumerate}
  \item What went well while writing this deliverable? 
  \item What pain points did you experience during this deliverable, and how
    did you resolve them?
  \item Which of your design decisions stemmed from speaking to your client(s)
  or a proxy (e.g. your peers, stakeholders, potential users)? For those that
  were not, why, and where did they come from?
  \item While creating the design doc, what parts of your other documents (e.g.
  requirements, hazard analysis, etc), it any, needed to be changed, and why?
  \item What are the limitations of your solution?  Put another way, given
  unlimited resources, what could you do to make the project better? (LO\_ProbSolutions)
  \item Give a brief overview of other design solutions you considered.  What
  are the benefits and tradeoffs of those other designs compared with the chosen
  design?  From all the potential options, why did you select the documented design?
  (LO\_Explores)
  \item What have you learned from implementing another team's module?
\end{enumerate}
\subsection*{Group - Reflection}
\begin{enumerate}
  \item N/A
  \item N/A
  \item N/A
  \item I think that our other documents such as the \href{https://github.com/M9Huynh/technically-functional/blob/main/docs/SRS/SRS.pdf}{requirements}, \href{https://github.com/M9Huynh/technically-functional/blob/main/docs/HazardAnalysis/HazardAnalysis.pdf}{hazard analysis} need to be slightly adjusted as a result of the design documents. With the requirements specifically some of them pertain to how the application design is and after completing our initial draft of the MG and MIS, the design does not necessarily line up with the requirements document. Further, some of the hazards considered should be changed to include a couple more areas for hazards to occur. Lastly for both of the documents, the assumptions should be changed to incorporated design changes that we had not considered during that deliverable. 
  \item Some limitations on the current system design and with the set scope include the accuracy of the pose estimation, the responsiveness of the application, and the limitation of stored exercises. With more resources, we would like to improve the accuracy of the pose estimation by either using a more advanced model or training our own model specifically for physiotherapy exercises. Further, we would like to improve the responsiveness of the application by optimizing our code and potentially using more powerful hardware to run the application. And finally, although explicity stated in our scope, we would like to include more exercises and potentially allow for custom exercises to be added by physiotherapists.
  
  \item One alternative design solution we considered was making the application a desktop application rather than an android application. The benefit of this design is that it would allow for more powerful hardware to be used, which could improve the responsiveness and accuracy of the application. However, the tradeoff is a limitation to the accessibility of the application, as not all users may have access to a desktop computer with hardware capable of recording. We ultimately chose the android application design because it allows for greater accessibility and convenience for users, as they can use the application on their mobile devices which are easier to setup and come with recording hardware.
  
  \item Implementing another team's module made it clear that another level of description is needed in order to effectively communicate the design and function of each module to developers. While implementing the Pixelation module for team 7, we found the description of their access program inputs and outputs lacking in detail. With sufficient context from the rest of the documentation we were able to produce what we believed to be the intended functionality of the module, but it required a lot of extra time and effort to interpret the design document. Additionally, because of this vague input and output description, there were no names given to input and output variables, which made it harder to stick to their coding conventions but also even harder to interpret the context of the module. 
  
  There was a lack of detail which I feel resonates with our documentation as well. The purposes and usage of each module by other modules should be clearly defined in order to avoid confusion during implementation. Working on the final version of the documentation, we will make sure to include detailed descriptions of each access program's inputs and outputs, as well as any necessary context for understanding the module's function within the overall system.

  There was an overall description of modules and their functionality which provided a good high-level overview of each module, something we think should be included in our documentation as well. Additionally the use of assumptions to clarify certain design decisions was helpful, and while we don't believe it to be the correct section to provide context it was a good signifier of a well made assumption.

  Overall, this experience has taught us the importance of clear and detailed documentation in order to facilitate effective communication and collaboration between teams.
\end{enumerate}

\subsection*{Matthew - Reflection}
\begin{enumerate}
  \item I think something that went well was our changes to overall workflow. One thing that we recognized was starting our document on google docs and then transferring our to sections to GitHub on the last two days the the deliverable was due led to some things going unchecked and the deliverable not turning out the way we had planned. Further it had also resulted in a 'panic' when we had to resolve merge conflicts between our sections soon to the deadline. Overall, we cut google docs for this deliverable and mainly focused on our sections within GitHub and doing multiple small PRs within the deliverable timeline and that had gone quite well. There are a few days left but a majority of the document has at least been accounted for and hopefully this workflow works for our team.
  \item I think one pain point I experienced was time management. I had suggested that since the computing team's workload was quite different to the documentation team's, we could take on small sections within the document with hopes of submitting it early into the project. This did partially go well but with focus split in two areas it was hard to focus on the development side of things for the POC demo. This may be something that improves with time but as of now is currently a work in progress as development on my end is not where I want it to be.
  \item I think that our design decisions partially came from speaking with our stakeholder (Kevin Yu) but other decisions were made with typical application creation. For example, the high-level was talked through with Kevin to simulate the timing and feedback delivery that you would encounter during physiotherapy but things like reporting, logging-in as well as account creation were from looking at other applications and aspects that went well from that end. 
  \item Group Q
  \item Group Q
  \item Group Q
  

\end{enumerate}

\subsection*{Vaisnavi - Reflection}
\begin{enumerate}
  \item Something that went well was how we worked on creating more constant pull requests (PRs) to avoid merge conflicts. This worked out well with everyone consistently pushing more PRs early on, allowing someone on the team to review and effectively merge it in. Overall, the organization and collaboration went very smoothly, and it allowed us to ultimately finish ahead of our internal deadline so we had enough time for editing and formatting.
  \item I think one of the major pain points through this deliverable was trying to balance both the development side for the POC demo while also trying to contribute to the document. As the main sub-team for the POC demo, we did take on smaller sections, but it was difficult to find that balance in time for both at the start. 
  \item We believe that most of our major design decisions stemmed directly from our meetings with our stakeholder, Kevin Yu. Throughout the development of this deliverable, a number of ideas that we originally planned had to be changed or completely re-worked based on the feedback we received. In several cases, Kevin helped us recognize that certain features were either not feasible within the course timeline or did not align with the true purpose and scope of our project. This was extremely helpful overall to have this input to tackle these decisions early on in the process.
  \item Group Q
  \item Group Q
  \item Group Q
  

\end{enumerate}
\subsection*{Maham - Reflection}
\begin{enumerate}
  \item Despite our (extremely) busy schedules, we made time to meet when we could and offer assistance to any member that needed it. In addition, all of us kept in contact and checked in with each other. I liked that none of us completely abandoned this project, but rather balanced it quite well. Also, the more frequent PR requests strategy (implemented for this deliverable) ended up being immensely helpful.
  \item A pain point for me was that a lot of my time and energy was taken up by external tasks. These were, unfortunately, appointments I could not compromise on because they had been scheduled months in advance. I would have loved to give this project more of my focus, and I tried my best to self-organize, but there were some areas I fell short. I will attempt to refine my balancing act for the future.
  \item The project is still in the early stages of development, and as a result, we do not have many concrete elements. However, some of developed ideas and design decisions was aided by Kevin, our supervisor. At this stage, we are constructing the foundation of the software, and that knowledge was supplied by him. Eventually, when we do have a functional product, we will want to loop in other stakeholders, e.g. patients, for modifications
  \item Group Q
  \item Group Q
  \item Group Q
\end{enumerate}

\subsection*{Cieran - Reflection}
\begin{enumerate}
  \item During this deliverable there was a need to agree upon implementation tactics. The team met and nailed down an approach for the development of the proof of concept, and I think everyone did a good job of ensuring their thoughts on what the system $was$ got at least mentioned when developing the system.
  \item I yet again sharked my teammates with procrastination. I have apologized profusely for my lack of initiative with this deliverable and have taken steps to ensure progress is far more consistent and meaningful in the future. As for this deliverable, I worked to catch up my own lack of effort towards the end of the it.
  \item I think the design decisions, specifically around the system's allowance for physiotherapists to alter patient's exercise routines, were very influenced by meetings with our advisor. As a team, we wanted to ensure that the system was something wanted by both patients and physiotherapists; designing based off of requirements from past physiotherapy patients and our physiotherapist advisor was a must. Some design decisions, such as checking against public databases for the physiotherapist's identification number, were made in an attempt at security for users of the application.
  \item GROUP Q
  \item GROUP Q
  \item GROUP Q
\end{enumerate}
\end{document}